\subsection{Decomposition and inertia for integrally closed domains}

LEMMA 4. Let A be an integrally closed domain, K its field of fractions, p the characteristics exponent of K, K' a radicial extension of K and A' a subring of K' containing A and integral over A. For every prime ideal p of A, there exists a unique prime ideal \(\mathfrak{p}'\) of \(\mathbf{A}'\) lying above \(\mathfrak{p}\) and \(\mathfrak{p}'\) is the set of \(x \in \mathbf{A}'\) such that there exists an integer \(m \geqslant 0\) for which \(x^{p^m} \in \mathfrak{p}\) .

The existence of \(\mathfrak{p}'\) follows from no. 1, Theorem 1. If \(x \in \mathfrak{p}'\) , there exists \(m \geqslant 0\) such that \(x^{p^m} \in \mathbf{K}\) , whence \(x^{p^m} \in \mathbf{A}\) since \(\mathbf{A}\) is integrally closed, hence \(x^{p^m} \in \mathfrak{p}' \cap \mathbf{A} = \mathfrak{p}\) . Conversely, if \(x \in \mathbf{A}'\) is such that \(x^{p^m} \in \mathfrak{p} \subset \mathfrak{p}'\) , then \(x \in \mathfrak{p}\) since \(\mathfrak{p}'\) is prime.

Remark (1). It follows from § 1, no. 3, Corollary to Proposition 11 that the integral closure of A in \(\mathbf{K}'\) is the set of \(x \in \mathbf{K}'\) such that there exists \(m \geqslant 0\) for which \(x^{p^m} \in \mathbf{A}\) (Algebra, Chapter V, § 8, no. 1, Proposition 1).

PROPOSITION 6. Let A be an integrally closed domain, K its field of fractions, K' a quasi-Galois extension of K and A' the integral closure of A in K'. Then:

\begin{enumerate}
\def\labelenumi{(\roman{enumi})}
\item
  For every prime ideal p of A, the group G of K-automorphisms of K' operates transitively on the set of prime ideals of A' lying over p.
\item
  For every prime ideal \(\mathfrak{p}'\) of \(\mathbf{A}'\) , the field of fractions \(k'\) of \(\mathbf{A}' / \mathfrak{p}'\) is a quasi-Galois extension of the field of fractions \(k\) of \(\mathbf{A} / (\mathbf{A} \cap \mathfrak{p}')\) and the canonical homomorphism \(\sigma \mapsto \bar{\sigma}\) from \(\mathcal{G}^{\mathbb{Z}}(\mathfrak{p}')\) to the group \(\Gamma\) of \(k\) -automorphisms of \(k'\) defines, by taking the quotient, a bijection of \(\mathcal{G}^{\mathbb{Z}}(\mathfrak{p}') / \mathcal{G}^{\mathbb{T}}(\mathfrak{p}')\) onto \(\Gamma\) .
\end{enumerate}

\begin{enumerate}
\def\labelenumi{(\Alph{enumi})}
\item
  Suppose first that \(\mathbf{K}'\) is a finite Galois extension of \(\mathbf{K}\) . Then \(\mathbf{A} = \mathbf{A}' \cap \mathbf{K}\) since \(\mathbf{A}\) is integrally closed and \(\mathbf{A}\) is therefore the ring of invariants of \(\mathcal{G}\) in \(\mathbf{A}'\) . As \(\mathcal{G}\) is finite, the proposition follows in this case from no. 2, Theorem 2.
\item
  Suppose secondly that \(\mathbf{K}'\) is any Galois extension of \(\mathbf{K}\) . Then \(\mathbf{K}'\) is the union of a right directed family \((\mathrm{K}_{\alpha})_{\alpha \in I}\) of finite Galois extensions of \(\mathbf{K}\) . To show (i), consider two prime ideals \(\mathfrak{p}'\) , \(\mathfrak{q}'\) of \(\mathbf{A}'\) lying above \(\mathfrak{p}\) . For all \(\alpha \in I\) , \(\mathfrak{p}' \cap \mathrm{K}_{\alpha}\) and \(\mathfrak{q}' \cap \mathrm{K}_{\alpha}\) are two prime ideals of \(\mathbf{A}' \cap \mathrm{K}_{\alpha}\) lying above \(\mathfrak{p}\) . Since \(\mathbf{A}' \cap \mathrm{K}_{\alpha}\) is the integral closure of \(A\) in \(\mathrm{K}_{\alpha}\) and the restrictions to \(\mathrm{K}_{\alpha}\) of the elements of \(\mathcal{G}\) form the group of K-automorphisms of \(\mathbf{K}\) , it follows from case (A) that there exists \(\sigma \in \mathcal{G}\) such that \(\sigma. (\mathfrak{p}' \cap \mathrm{K}_{\alpha}) = \mathfrak{q}' \cap \mathrm{K}_{\alpha}\) . Let \(\mathcal{E}_{\alpha}\) be the set of \(\sigma \in \mathcal{G}\) which have the above property. Let \(\sigma \in \mathcal{G} - \mathcal{E}_{\alpha}\) ; then, for all \(\tau \in \mathcal{G}\) leaving invariant the elements of \(\mathrm{K}_{\alpha}, (\sigma \tau). (\mathfrak{p}' \cap \mathrm{K}_{\alpha}) = \sigma. (\mathfrak{p}' \cap \mathrm{K}_{\alpha}) \neq \mathfrak{q}' \cap \mathrm{K}_{\alpha}\) and hence \(\sigma \tau \in \mathcal{G} - \mathcal{E}_{\alpha}\) . It follows that \(\mathcal{E}_{\alpha}\) is closed in the topological Galois group \(\mathcal{G}\) (Algebra, Chapter V, Appendix II, no. 1) and clearly the family \((\mathcal{E}_{\alpha})_{\alpha \in I}\) is left directed. As \(\mathcal{G}\) is compact (loc. cit., no. 2, Proposition 3) and the \(\mathcal{E}_{\alpha}\) are non-empty, the intersection \(\mathcal{E}\) of the family \((\mathcal{E}_{\alpha})\) is non-empty and \(\sigma. \mathfrak{p}' = \mathfrak{q}'\) for all \(\sigma \in \mathcal{E}\) , whence (i).
\end{enumerate}

To show (ii), note that \(k'\) is the union of the right directed family \((k_{\alpha})_{\alpha \in \mathbf{I}}\) , where \(k_{\alpha}\) is the field of fractions of \((\mathrm{A}' \cap \mathrm{K}_{\alpha}) / (\mathfrak{p}' \cap \mathrm{K}_{\alpha})\) . As each \(k_{\alpha}\) is a quasi-Galois extension of \(k\) by (A), so is \(k'\) (Algebra, Chapter V, § 6, no. 3, Proposition 8). On the other hand, let \(u\) be a \(k\) -automorphism of \(k'\) and let \(\pi': \mathrm{A}' \to \mathrm{A}' / \mathfrak{p}'\) be the canonical homomorphism. By virtue of no. 2, Theorem 2 applied to \(\mathrm{A}' \cap \mathrm{K}_{\alpha}\) , there exists for all \(\alpha\) a non-empty set \(\mathcal{F}_{\alpha}\) of elements \(\sigma \in \mathcal{G}\) such that \(\sigma \cdot (\mathfrak{p}' \cap \mathrm{K}_{\alpha}) = \mathfrak{p}' \cap \mathrm{K}_{\alpha}\) and \(u(\pi'(x)) = \pi'(\sigma.x)\) for all \(x \in \mathrm{A}' \cap \mathrm{K}_{\alpha}\) . As above it is seen that \(\mathcal{F}_{\alpha}\) is closed in \(\mathcal{G}\) and, as \((\mathcal{F}_{\alpha})\) is a left directed set, its intersection \(\mathcal{F}\) is non-empty. Clearly for \(\sigma \in \mathcal{F}\) , \(\sigma \in \mathcal{G}^{\mathbb{Z}}(\mathfrak{p}')\) and \(\bar{\sigma} = u\) , which completes the proof of (ii) in this case.

\begin{enumerate}
\def\labelenumi{(\Alph{enumi})}
\setcounter{enumi}{2}
\tightlist
\item
  General case. The field of invariants \(K_{1}\) of G is a radicial extension of K (Algebra, Chapter V, § 10, no. 9, Proposition 14); there therefore exists a single prime ideal \(p_{1}\) of \(A_{1} = A' \cap K_{1}\) lying above p (Lemma 4). If \(p'\) and \(q'\) are two prime ideals of \(A'\) lying above p, then they lie above \(p_{1}\) ; as \(K'\) is a Galois extension of \(K_{1}\) and \(A' \cap K_{1}\) is integrally closed (§ 1, no. 2, Proposition 7 and Corollary to Proposition 8), it follows from (B) that there exists \(\sigma \in G\) such that \(\sigma \cdot p' = q'\) ; whence (i). On the other hand, clearly the field of fractions \(k_{1}\) of \(A_{1}/p_{1}\) is a radicial extension of k (A being integrally closed); as \(k'\) is a quasi-Galois extension of \(k_{1}\) by (B), \(k'\) is a quasi-Galois extension of k, every k-isomorphism of \(k'\) to an algebraically closed extension of \(k'\) being a \(k_{1}\) -isomorphism. This last remark shows also, taking account of (B), that every k-automorphism of \(k'\) is of the form \(\bar{\sigma}\) where \(\sigma \in \mathcal{G}^{\mathbb{Z}}(\mathfrak{p}')\) , which completes the proof of (ii).
\end{enumerate}

Remark

\begin{enumerate}
\def\labelenumi{(\arabic{enumi})}
\setcounter{enumi}{1}
\tightlist
\item
  Suppose that \(\mathbf{K}'\) is a Galois extension of \(\mathbf{K}\) and let us keep the notation of the proof of Proposition 6; for all \(\alpha\) let \(\mathcal{G}_{\alpha}^{\mathbb{Z}}\) (resp. \(\mathcal{G}_{\alpha}^{\mathrm{T}}\) ) be the subgroup of \(\mathcal{G}\) consisting of the \(\sigma\) whose restriction to \(\mathrm{A}' \bigcap \mathrm{K}_{\alpha}\) belongs to \(\mathcal{G}^{\mathbb{Z}}(\mathfrak{p}' \bigcap \mathrm{K}_{\alpha})\) (resp. \(\mathcal{G}^{\mathrm{T}}(\mathfrak{p}' \bigcap \mathrm{K}_{\alpha})\) ). The proof of Proposition 6 shows that these subgroups are closed in \(\mathcal{G}\) and that
\end{enumerate}

\[
\mathcal {G} ^ {\mathrm{Z}} \left(\mathfrak {p} ^ {\prime}\right) = \bigcap_ {\alpha} \mathcal {G} _ {\alpha} ^ {\mathrm{Z}} \quad \text { and } \quad \mathcal {G} ^ {\mathrm{T}} \left(\mathfrak {p} ^ {\prime}\right) = \bigcap_ {\alpha} \mathcal {G} _ {\alpha} ^ {\mathrm{T}}.
\]

Moreover, the set of restrictions to \(A^{\prime} \cap K_{\alpha}\) of the elements of G (resp. \(G_{\alpha}^{T}\) ) is the whole of the group \(\mathcal{G}^{\mathtt{Z}}(\mathfrak{p}^{\prime} \cap \mathtt{K}_{\alpha})\) (resp. \(\mathcal{G}^{\mathtt{T}}(\mathfrak{p}^{\prime} \cap \mathtt{K}_{\alpha})\) ), every K-automorphism of \(K_{\alpha}\) extending to an element of G.

With the same hypotheses, the ring \(\mathbf{A}^{\mathbf{Z}}(\mathfrak{p}')\) (resp. \(\mathbf{A}^{\mathbf{T}}(\mathfrak{p}')\) ) is the union of the directed family of the \(\mathbf{A}^{\mathbf{Z}}(\mathfrak{p}' \cap \mathbf{K}_{\alpha})\) (resp. \(\mathbf{A}^{\mathbf{T}}(\mathfrak{p}' \cap \mathbf{K}_{\alpha})\) ): in fact, every \(x \in \mathbf{A}^{\mathbf{Z}}(\mathfrak{p}')\) (resp. every \(x \in \mathbf{A}^{\mathbf{T}}(\mathfrak{p}')\) ) belongs to one of the \(\mathbf{K}_{\alpha}\) and by the above there exists a \(\beta\) such that \(\mathbf{K}_{\alpha} \subset \mathbf{K}_{\beta}\) and the restrictions to \(\mathbf{A}' \cap \mathbf{K}_{\alpha}\) of the elements of \(\mathcal{G}^{\mathbf{Z}}(\mathfrak{p}')\) (resp. \(\mathcal{G}^{\mathbf{T}}(\mathfrak{p}')\) ) are the same as the restrictions to \(\mathbf{A}' \cap \mathbf{K}_{\alpha}\) of the elements of \(\mathcal{G}^{\mathbf{Z}}(\mathfrak{p}' \cap \mathbf{K}_{\beta})\) (resp. \(\mathcal{G}^{\mathbf{T}}(\mathfrak{p}' \cap \mathbf{K}_{\beta})\) ), the groups \(\mathcal{G}^{\mathbf{Z}}(\mathfrak{p}' \cap \mathbf{K}_{\alpha})\) and \(\mathcal{G}^{\mathbf{T}}(\mathfrak{p}' \cap \mathbf{K}_{\beta})\) being finite; hence \(x\) belongs to \(\mathbf{A}^{\mathbf{Z}}(\mathfrak{p}' \cap \mathbf{K}_{\beta})\) (resp. \(\mathbf{A}^{\mathbf{T}}(\mathfrak{p}' \cap \mathbf{K}_{\beta})\) ).

COROLLARY 1. Under the hypotheses of Proposition 6, let \(f\) be a homomorphism from \(\mathbf{A}\) to a field \(\mathbf{L}\) and \(g_1\) , \(g_2\) two homomorphisms from \(\mathbf{A}'\) to \(\mathbf{L}\) which extend \(f\) . Then there exists a \(\mathbf{K}\) -automorphism \(\sigma\) of \(\mathbf{K}'\) such that \(g_1 = g_2 \circ \sigma\) .

The proof starting from Proposition 6 is the same as that of the Corollary to Theorem 2 starting from the latter.

COROLLARY 2. Let A be an integrally closed domain, K its field of fractions, K' a finite algebraic extension of K and A' a subring of K' containing A and integral over A.

\begin{enumerate}
\def\labelenumi{(\roman{enumi})}
\item
  For every prime ideal \(\mathfrak{p}\) of \(\mathbf{A}\) the set of prime ideals of \(\mathbf{A}'\) lying above \(\mathfrak{p}\) is finite.\\
\item
  If \(\mathfrak{p}'\) is a prime ideal of \(\mathbf{A}'\) lying above \(\mathfrak{p}\) , every element of \(\mathbf{A}' / \mathfrak{p}'\) is of degree \(\leqslant [\mathbf{K}' : \mathbf{K}]\) over the field of fractions of \(\mathbf{A} / \mathfrak{p}\) .
\item
  Let \(\mathbf{K}''\) be the quasi-Galois extension of \(\mathbf{K}\) generated by \(\mathbf{K}'\) in an algebraic closure of \(\mathbf{K}'\) and \(\mathbf{A}''\) the integral closure of \(\mathbf{A}\) in \(\mathbf{K}''\) . The field \(\mathbf{K}''\) is a finite extension of \(\mathbf{K}\) (Algebra, Chapter V, § 6, no. 3, Corollary 1 to Proposition 9) and hence its group of K-automorphisms is finite; it follows that the set of prime ideals of \(\mathbf{A}''\) lying above \(p\) is finite (Proposition 6 (i)). On the other hand, as \(\mathbf{A}''\) is integral over \(\mathbf{A}\) , the mapping \(p'' \mapsto p'' \cap \mathbf{A}'\) of the set of prime ideals of \(\mathbf{A}''\) lying above \(p\) to the set of prime ideals of \(\mathbf{A}'\) lying above \(x\) is surjective (no. 1, Theorem 1).
\item
  The coefficients of the minimal polynomial (over K) of any element \(x' \in A'\) belong to A (§ 1, no. 3, Corollary to Proposition 10); applying the canonical homomorphism \(\pi': A' \to A'/p'\) to the coefficients of this polynomial, an equation of integral dependence with coefficients in A/p and of degree \(\leqslant [K': K]\) is obtained for the class mod. \(p'\) of \(x\) ; whence the conclusion.
\end{enumerate}

COROLLARY 3. With the hypotheses and notation of Corollary 2, if A is semi-local, so is A'.

For every maximal ideal \(m'\) of \(A'\) , \(m' \cap A\) is a maximal ideal of A (no. 1, Proposition 1); the corollary then follows from Corollary 2, since by hypothesis the set of maximal ideals of A is finite.

COROLLARY 4. Let A be an integrally closed domain, K its field of fractions, K' a Galois extension of K, A' the integral closure of A in K', p' a prime ideal of A', p = A ∩ p' and k and k' the fields of fractions of A/p and A'/p' respectively. Then:

\begin{enumerate}
\def\labelenumi{(\roman{enumi})}
\item
  The field of fractions of \(\mathbf{A}^{\mathbf{z}} / (\mathfrak{p}'\cap \mathbf{A}^{\mathbf{z}})\) is equal to \(k\) and the maximal ideal of the local ring of \(\mathbf{A}^{\mathbf{z}}\) relative to \(\mathfrak{p}'\cap \mathbf{A}^{\mathbf{z}}\) is generated by \(\mathfrak{p}\) .
\item
  The field of fractions \(k^{\mathbf{T}}\) of \(\mathbf{A}^{\mathbf{T}} / (\mathfrak{p}' \cap \mathbf{A}^{\mathbf{T}})\) is the greatest separable extension of \(k\) contained in \(k'\) .
\end{enumerate}

The ring A is the ring of invariants in A' of the Galois group of K' over K; if K' is of finite degree over K, the corollary then follows from Propositions 4 and 5 of no. 2. Consider now the general case, K' therefore being the union of a right directed set (K \(_{\alpha}\) ) of finite Galois extensions of K. Then:

\begin{enumerate}
\def\labelenumi{(\roman{enumi})}
\tightlist
\item
  If \(x, y\) are two elements of \(\mathbf{A}^z\) , where \(y \notin \mathfrak{p}'\) , there is an index \(\alpha\) such that \(x\) and \(y\) belong to \(\mathbf{A}^z(\mathfrak{p}' \cap \mathbf{K}_\alpha)\) (Remark 2); by Proposition 4 of no. 2, there are \(x_0, y_0\) in \(\mathbf{A}\) with \(y_0 \notin \mathfrak{p}'\) such that \(xy_0 - x_0y \in \mathfrak{p}'\) , which proves the first assertion of (i); if further \(x \in \mathfrak{p}'\) , we may assume that \(y_0\) satisfies
\end{enumerate}

\[
x y _ {0} \in \mathfrak {p A} ^ {\mathrm{z}} (\mathfrak {p} ^ {\prime} \cap \mathrm{K} _ {\alpha}) \subset \mathfrak {p A} ^ {\mathrm{z}} (\mathfrak {p} ^ {\prime}),
\]

which proves the second assertion of (i).

\begin{enumerate}
\def\labelenumi{(\roman{enumi})}
\setcounter{enumi}{1}
\tightlist
\item
  Suppose now that \(x \in \mathbf{A}^{\mathrm{T}}\) ; there exists \(\alpha\) such that \(x \in \mathbf{A}^{\mathrm{T}}(\mathfrak{p}' \cap \mathbf{K}_{\alpha})\) (Remark 2) and Proposition 5 of no. 2 shows that the class \(\bar{x}\) of \(x \bmod. (\mathfrak{p}' \cap \mathbf{K}_{\alpha} \cap \mathbf{A}^{\mathrm{T}})\) is algebraic and separable over \(k\) ; a fortiori the class mod. \((\mathfrak{p}' \cap \mathbf{A}^{\mathrm{T}})\) of \(x\) is separable over \(k\) ; to complete the proof of the corollary, it is sufficient to show that \(k'\) is a radicial extension of \(k^{\mathrm{T}}\) . Now, \(k'\) is the union of the right directed family of fields of fractions \(k_{\alpha}\) of the rings \((\mathbf{A}' \cap \mathbf{K}_{\alpha}) / (\mathfrak{p}' \cap \mathbf{K}_{\alpha})\) . It follows therefore from Proposition 5 that, if an element of \(k'\) belongs to \(k_{\alpha}\) , it is radicial over the field of fractions of
\end{enumerate}

\[
\mathbf {A} ^ {\mathrm{T}} \left(\mathfrak {p} ^ {\prime} \cap \mathbf {K} _ {\alpha}\right) / \left(\mathfrak {p} ^ {\prime} \cap \mathbf {A} ^ {\mathrm{T}} \left(\mathfrak {p} ^ {\prime} \cap \mathbf {K} _ {\alpha}\right)\right)
\]

and a fortiori over \(k^{\mathbf{T}}\) (by virtue of Remark 2).

DEFINITION 5. With the hypotheses and notation of Proposition 6, the field of invariants \(\mathbf{K}^{\mathrm{Z}}(\mathfrak{p}')\) (resp. \(\mathbf{K}^{\mathrm{T}}(\mathfrak{p}')\) ) of the group \(\mathcal{G}^{\mathrm{Z}}(\mathfrak{p}')\) (resp. \(\mathcal{G}^{\mathrm{T}}(\mathfrak{p}')\) ) in the field \(\mathbf{K}'\) is called the decomposition field (resp. inertia field) of \(\mathfrak{p}'\) with respect to \(\mathbf{K}\) .

We also write \(K^{z}\) (resp. \(K^{T}\) ) in place of \(\mathbf{K}^{\mathbf{Z}}(\mathfrak{p}')\) (resp. \(\mathbf{K}^{\mathrm{T}}(\mathfrak{p}')\) ). It follows from §1, no. 9, Proposition 23 that \(K^{z}\) (resp. \(K^{T}\) ) is the field of fractions of the ring \(A^{z}\) (resp. \(A^{T}\) ); \(A^{z}\) (resp. \(A^{T}\) ) is the integral closure of A in \(K^{z}\) (resp. \(K^{T}\) ).

\textbf{Remarks}\par

\begin{enumerate}
\def\labelenumi{(\arabic{enumi})}
\setcounter{enumi}{2}
\tightlist
\item
  Under the conditions of Corollary 4 of Proposition 6 and assuming that \([\mathbf{K}':\mathbf{K}]\) is finite, the number of distinct prime ideals lying above \(p\) is \([\mathbf{K}^{\mathbf{Z}}:\mathbf{K}]\) , this degree being equal to the index \((\mathcal{G}:\mathcal{G}^{\mathbf{Z}})\) by Galois theory; moreover, it follows from Galois theory that
\end{enumerate}

\[
[ \mathbf {K} ^ {\mathrm{T}}: \mathbf {K} ^ {\mathrm{Z}} ] = (\mathcal {G} ^ {\mathrm{Z}}: \mathcal {G} ^ {\mathrm{T}}) = [ k ^ {\mathrm{T}}: k ].\tag{6}
\]

\begin{enumerate}
\def\labelenumi{(\arabic{enumi})}
\setcounter{enumi}{3}
\tightlist
\item
  Let A be an integrally closed domain, K its field of fractions, \(K'\) a finite algebraic extension of K and \(A'\) the integral closure of A in \(K'\) . Then, for every prime ideal p of A, the number of prime ideals of \(A'\) lying above p is at most \([K':K]_{s}\) (the separable factor of the degree of \(K'\) over K). We may first restrict our attention to the case where \(K'\) is a separable extension of K, for in general \(K'\) is a radicial extension of the greatest separable extension \(K_{0}\) of K contained in \(K'\) , \([K':K]_{s} = [K_{0}:K]\) by definition and, if \(A_{0}\) is the integral closure of A in \(K_{0}\) , the prime ideals of \(A_{0}\) and \(A'\) are in one-to-one correspondence (Lemma 4). Suppose therefore that \(K'\) is separable over K and let N be the Galois extension of K generated by \(K'\) in an algebraic closure of K, G its Galois group, B the integral closure of A in N and P a prime ideal of B lying above p.~Let H be the Galois group of N over \(K'\) and \(G^{z}\) the decomposition group of P; the prime ideals of B lying above p are the s.P where \(s \in G\) (no. 2, Theorem 2) and the relation s.P = s'.P means that \(s' = sg\) where \(g \in G^{z}\) . On the other hand in order that s.P ∩ \(K' = s'\) .P ∩ \(K'\) , it is necessary and sufficient that \(s'\) .P = ts.P, where \(t \in H\) (no. 2, Theorem 2), whence finally \(s' = tsg\) where \(t \in H\) and \(g \in G^{z}\) . The number of prime ideals of A' lying above p is therefore equal to the number of classes of G under the equivalence relation ``there exists \(t \in H\) and \(g \in G^{z}\) such that \(s' = tsg\)'' between s and \(s'\) ; clearly this number is at most equal to the index \((\mathcal{G}:\mathcal{H})\) , the number of right cosets of H in G, and \((\mathcal{G}:\mathcal{H}) = [\mathrm{K}':\mathrm{K}]\) by Galois theory.
\end{enumerate}

PROPOSITION 7. Let A be an integrally closed domain, K its field of fractions, K' a Galois extension of K, G its Galois group, A' the integral closure of A in K', p' a prime ideal of A' and p = A ∩ p'. Finally, let L be a subfield of K' containing K and let p(L) = p' ∩ L.

\begin{enumerate}
\def\labelenumi{(\roman{enumi})}
\item
  The decomposition field (resp. inertia field) of \(\mathfrak{p}'\) with respect to \(\mathbf{L}\) is \(\mathbf{L}(\mathbf{K}^{\mathbf{Z}})\) (resp. \(\mathbf{L}(\mathbf{K}^{\mathbf{T}})\) ); if further \(\mathbf{L}\) is a Galois extension of \(\mathbf{K}\) , the decomposition field of \(\mathfrak{p}(\mathbf{L})\) with respect to \(\mathbf{K}\) is \(\mathbf{L} \cap \mathbf{K}^{\mathbf{Z}}\) .
\item
  If \(\mathbf{L}\) is contained in \(\mathbf{K}^{\mathbf{Z}}\) , \(\mathbf{A} / \mathfrak{p}\) and \((\mathbf{A}' \cap \mathbf{L}) / \mathfrak{p}(\mathbf{L})\) have the same field of fractions and in the local ring \(\mathbf{A}' \cap \mathbf{L}\) corresponding to the prime ideal \(\mathfrak{p}(\mathbf{L})\) , the maximal ideal is generated by \(\mathfrak{p}\) . Conversely, if these two conditions hold and further \(\bigcap_{n \geqslant 0} \mathfrak{p}^{n} \mathbf{A}_{\mathfrak{p}'}' = 0\) , \(\mathbf{L}\) is contained in \(\mathbf{K}^{\mathbf{Z}}\) .
\item
  If \(\mathcal{H}\) is the subgroup of \(\mathcal{G}\) leaving invariant the elements of \(\mathbf{L}\) , clearly the decomposition group (resp. inertia group) of \(\mathfrak{p}'\) with respect to \(\mathbf{L}\) is \(\mathcal{G}^{\mathrm{Z}}\cap \mathcal{H}\) (resp. \(\mathcal{G}^{\mathrm{T}}\cap \mathcal{H}\) ) and the first assertion follows from Galois theory if \(\mathbf{K}'\) is a finite Galois extension of \(\mathbf{K}\) (Algebra, Chapter V, § 10, no. 6, Corollary 1 to Theorem 3); in the general case it follows from the fact that \(\mathbf{A}^{\mathrm{Z}}\) (resp. \(\mathbf{A}^{\mathrm{T}}\) ) is the union of the \(\mathbf{A}^{\mathrm{Z}}(\mathfrak{p}'\cap \mathbf{K}_{\alpha})\) (resp. \(\mathbf{A}^{\mathrm{T}}(\mathfrak{p}'\cap \mathbf{K}_{\alpha})\) ) in the notation of Remark 2: every element \(x\in \mathbf{K}'\) belongs to some \(\mathbf{K}_{\alpha}\) and if it is invariant under \(\mathcal{G}^{\mathrm{Z}}(\mathfrak{p}')\cap \mathcal{H}\) (resp. \(\mathcal{G}^{\mathrm{T}}(\mathfrak{p}')\cap \mathcal{H}\) ) it is also invariant under \(\mathcal{G}^{\mathrm{Z}}(\mathfrak{p}'\cap \mathbf{K}_{\beta})\cap \mathcal{H}\) (resp. \(\mathcal{G}^{\mathrm{T}}(\mathfrak{p}'\cap \mathbf{K}_{\beta})\cap \mathcal{H}\) ) for some suitable \(\beta\) ; hence it belongs by the beginning of the argument to
\end{enumerate}

\[
\mathbf {L} \left(\mathrm{K} ^ {\mathrm{Z}} \left(\mathfrak {p} ^ {\prime} \cap \mathrm{K} _ {\alpha}\right)\right) \subset \mathbf {L} \left(\mathrm{K} ^ {\mathrm{Z}}\right) (\text { resp. } \mathbf {L} \left(\mathrm{K} ^ {\mathrm{T}} \left(\mathfrak {p} ^ {\prime} \cap \mathrm{K} _ {\beta}\right)\right) \subset \mathbf {L} \left(\mathrm{K} ^ {\mathrm{T}}\right)).
\]

Suppose now that L is a Galois extension of K; the restriction to L of every \(\sigma\in G^{z}\) then leaves invariant \(\mathfrak{p}(L)=\mathfrak{p}^{\prime}\cap L\) and hence belongs to the decomposition group of \(\mathfrak{p}(L)\) with respect to K. Conversely, let \(\tau\) be an automorphism of L belonging to this group and let \(\sigma\) be an extension of \(\tau\) to a K-automorphism of \(K^{\prime}\) ; we write \(q^{\prime}=\sigma.p^{\prime}\) . As \(p^{\prime}\) and \(q^{\prime}\) both lie above \(\mathfrak{p}(L)\) , there exists an automorphism \(\rho\in H\) such that \(q^{\prime}=\rho.p^{\prime}\) , whence \(\rho^{-1}\sigma\in G^{z}\) and \(\tau\) is the restriction of \(\rho^{-1}\sigma\) to L; in other words, the decomposition group of \(\mathfrak{p}(L)\) with respect to K is identical with the group of restrictions to L of the automorphisms \(\sigma\in G^{z}\) , which proves the second assertion.

\begin{enumerate}
\def\labelenumi{(\roman{enumi})}
\setcounter{enumi}{1}
\tightlist
\item
  To say that \(L \subset K^{z}\) means that \(H \supset G^{z}\) and the assertions of (ii) are therefore special cases of no. 2, Proposition 4 (ii) and (iii) when \([K':K]\) is finite. In the general case the argument is as in the proof of Proposition 6.
\end{enumerate}

\subsection{The second existence theorem}

THEOREM 3. Let A be an integrally closed domain and A' a ring containing A and integral over A. Suppose that 0 is the only element of A which is a divisor of 0 in A'. Let p, q be two prime ideals of A such that q ⊃ p and q' a prime ideal of A' lying above q. Then there exists a prime ideal p' of A' lying above p and such that q' ⊃ p'.

Suppose first that A' is an integral domain. Let K, K' be the fields of fractions of A and A' respectively; let K'' be the algebraic closure of K' and A'' the integral closure of A in K''; then \(A \subset A' \subset A''\) . Let \(p''\) be a prime ideal of A'' lying above p (no. 1, Theorem 1), \(q''\) a prime ideal of A'' lying above q and such that \(p'' \subset q''\) (no. 1, Corollary 2 to Theorem 1) and finally \(q_{1}''\) a prime ideal of A'' lying above \(q'\) (no. 1, Theorem 1). By no. 3, Proposition 6 (i), there exists a K-automorphism \(\sigma\) of K'' such that \(\sigma \cdot q'' = q_{1}''\) . Then \(\sigma \cdot p''\) is a prime ideal of A'' lying above p such that \(\sigma \cdot p'' \subset q_{1}''\) and hence \(p' = A' \cap \sigma \cdot p''\) is a prime ideal of A' lying above p and contained in \(A' \cap q_{1}'' = q'\) .

We pass to the general case. As A is an integral domain and \(q'\) is prime, the subsets \(A - \{0\}\) and \(A' - q'\) of \(A'\) are multiplicative; then their product \(S = (A - \{0\})(A' - q')\) is a multiplicative subset of \(A'\) which does not contain 0 since the non-zero elements of A are not divisors of 0 in \(A'\) . Then there exists (Chapter II, §2, no. 5, Corollary 2 to Proposition 11) a prime ideal \(m'\) of \(A'\) disjoint from S, in other words such that \(m' \subset q'\) and \(m' \cap A = 0\) . Let h be the canonical homomorphism \(A' \to A'/m'\) . The restriction of h to A is injective and hence \(h(A)\) is integrally closed. As \(m' \subset q'\) , \(h(q')\) is a prime ideal of \(A'/m'\) lying above \(h(q)\) ; since \(A'/m'\) is an integral domain, the first part of the proof proves that there exists a prime ideal \(n'\) of \(A'/m'\) such that \(n' \cap h(A) = h(p)\) and \(h(q') \supset n'\) . The ideal \(p' = h^{-1}(n')\) is a prime ideal of \(A'\) and \(q' \supset p'\) , since \(q'\) contains the kernel of h. As \(n' \supset h(p)\) , \(p' \supset p\) . Finally, for \(x \in p' \cap A\) , \(h(x) \in n' \cap h(A) = h(p)\) and hence \(x \in p\) since the restriction of h to A is injective; hence \(p' \cap A = p\) .

COROLLARY. Under the hypotheses on A and A' of Theorem 3, let p be a prime ideal of A. The prime ideals of A' lying above p are the minimal elements of the set E of prime ideals of A' containing pA'.

A prime ideal of A' lying above p is minimal in E by virtue of no. 1, Corollary 1 to Proposition 1. Conversely, let q' be a minimal element of E. As q' ∩ A ⊃ p, Theorem 3 shows that there exists a prime ideal p' of A' lying above p such that q' ⊃ p'. As p' contains pA', the hypothesis made on q' implies that q' = p' and hence q' lies above p.

* Let V, V' be two affine algebraic varieties and f a morphism of V' to V such that \(f(V')\) is dense in V. Let A (resp. A') be the ring of functions regular on V (resp. V'); having been given f we can identify A with a subring of A'; suppose that A' is integral over A. Theorem 1 of no. 1 shows that for every irreducible subvariety W of V there exists an irreducible subvariety \(W'\) of \(V'\) such that \(f(W')\) is a dense subset of W; in particular, every point of V is the image of an irreducible subvariety of \(V'\) , which shows that f is surjective. Similarly, the restriction of f to every irreducible subvariety \(W'\) of \(V'\) maps \(W'\) onto an irreducible subvariety of V. Corollary 2 to Theorem 1, no. 1 shows that, if W and X are two irreducible subvarieties of V such that \(W \supset X\) and if \(W'\) is an irreducible subvariety of \(V'\) such that \(f(W') = W\) , then there exists an irreducible subvariety \(X'\) of \(V'\) contained in \(W'\) and such that \(f(X') = X\) .

If A is integrally closed, V is called normal. Theorem 3 shows that, if V is normal, if W and X are irreducible subvarieties of V such that \(W \supset X\) and if \(X'\) is an irreducible subvariety of \(V'\) such that \(f(\mathbf{X}') = \mathbf{X}\) , then there exists a subvariety \(W'\) of \(V'\) containing \(X'\) and such that \(f(W') = W\) . Finally, the Corollary to Theorem 3 shows that, if V is normal and W is an irreducible subvariety of V, the irreducible subvarieties \(W'\) of \(V'\) such that \(f(W') = W\) are just the irreducible components of \(f(W)\) .

\section{FINITELY GENERATED ALGEBRAS OVER A FIELD}

\subsection{The normalization lemma}

In this no. and the following, k denotes a commutative field.

THEOREM 1 (Normalization Lemma). Let A be a finitely generated \(k\) -algebra and let \(\mathfrak{a}_1 \subset \mathfrak{a}_2 \subset \ldots \subset \mathfrak{a}_p\) be an increasing finite sequence of ideals of A such that \(p \geqslant 1\) and \(\mathfrak{a}_p \neq \mathbf{A}\) . There exists a finite sequence \((x_i)_{1 \leqslant i \leqslant n}\) of elements of A which are algebraically independent over \(k\) (Chapter III, § 1, no. 1) and such that:

\begin{enumerate}
\def\labelenumi{(\alph{enumi})}
\item
  A is integral over the ring \(\mathbf{B} = k[x_1, \ldots, x_n]\) .
\item
  For all \(j\) such that \(1 \leqslant j \leqslant p\) , there exists an increasing sequence \((h(j))_{1 \leqslant j \leqslant p}\) of integers such that for all \(j\) the ideal \(\mathfrak{a}_j \cap \mathbf{B}\) of \(\mathbf{B}\) is generated \(x_1, \ldots, x_{h(j)}\) .
\end{enumerate}

Note first that it is sufficient to prove the theorem when A is a polynomial algebra \(k[Y_{1},\ldots,Y_{m}]\) . For in the general case, A is isomorphic to a quotient of such an algebra \(A'\) by an ideal \(a_{0}'\) ; let \(a_{j}'\) denote the inverse image of \(a_{j}\) in \(A'\) and let \(x_{i}'(1\leqslant i\leqslant r)\) be elements of \(A'\) satisfying the conditions of the statement for the ring \(A'\) and the increasing sequence of ideals \(a_{0}'\subset a_{1}'\subset\cdots\subset a_{p}'\) . Then the images \(x_{i}\) of the \(x_{i}'\) in A for \(i>h(0)\) satisfy the desired conditions; this is obvious for condition (b) and for condition (a) this follows from §1, no. 1, Proposition 2; finally, if the \(x_{i}(h(0)+1\leqslant i\leqslant r)\) were not algebraically independent over k, there would be a non-zero polynomial

\[
\mathrm{Q} \in k [ \mathrm{X} _ {h (0) + 1}, \dots , \mathrm{X} _ {r} ]
\]

such that \(\mathrm{Q}(x_{h(0)+1}^{\prime},\ldots,x_{r}^{\prime})\in\mathfrak{a}_{0}^{\prime}\cap\mathbf{B}^{\prime}\) , writing \(B^{\prime}=k[x_{1}^{\prime},\ldots,x_{r}^{\prime}]\) ; but by hypothesis, every element of \(a_{0}^{\prime}\cap B^{\prime}\) may be written uniquely as a polynomial in the \(x_{j}^{\prime}\ (1\leqslant j\leqslant r)\) with coefficients in k, each monomial of which contains at least one of the \(x_{j}^{\prime}\) for \(1\leqslant j\leqslant h(0)\) ; we therefore obtain a contradiction, which proves our assertion.

We shall therefore suppose in the rest of the proof that \(A = k[Y_{1}, \ldots, Y_{m}]\) and we shall argue by induction on p.

\begin{enumerate}
\def\labelenumi{(\Alph{enumi})}
\tightlist
\item
  p = 1. We distinguish two cases:
\end{enumerate}

\textbf{(A1) The ideal \(\mathfrak{a}_1\) is a principal ideal generated by an element \(x_1 \notin k\) .}\par

\(x_{1} = \mathrm{P}(\mathrm{Y}_{1},\ldots ,\mathrm{Y}_{m})\) , where \(\mathbf{P}\) is a non-constant polynomial. We shall see that, for a suitable choice of integers \(r_i > 0\) , the ring A is integral over \(\mathbf{B} = k[x_1,x_2,\dots ,x_m]\) , where

\[
x _ {i} = \mathrm{Y} _ {i} - \mathrm{Y} _ {1} ^ {r _ {i}} \quad (2 \leqslant i \leqslant m).\tag{1}
\]

For this it is sufficient to choose the \(r_i\) such that \(Y_1\) is integral over \(B\) (§ 1, no. 1, Proposition 4). Now, there is the relation

\[
\mathrm{P} \left(\mathrm{Y} _ {1}, x _ {2} + \mathrm{Y} _ {1} ^ {r _ {2}}, \dots , x _ {m} + \mathrm{Y} _ {1} ^ {r _ {m}}\right) - x _ {1} = 0.\tag{2}
\]

Let \(\mathbf{P} = \sum_{\mathbf{p}} a_{\mathbf{p}} \mathbf{Y}^{\mathbf{p}}\) , where \(\mathbf{p} = (p_1, \ldots, p_m)\) , the \(p_i\) being integers \(\geqslant 0\) , \(\mathbf{Y}^{\mathbf{p}} = \mathbf{Y}_1^{p_1} \ldots \mathbf{Y}_m^{p_m}\) , the \(a_{\mathbf{p}}\) are elements \(\neq 0\) of \(k\) and at least one of the indices \(\mathbf{p}\) is distinct from \((0, \ldots, 0)\) ; relation (2) may be written

\[
\sum_ {\mathbf {p}} a _ {\mathbf {p}} \mathrm{Y} _ {1} ^ {p _ {1}} (x _ {2} + \mathrm{Y} _ {1} ^ {r _ {2}}) ^ {p _ {2}} \dots (x _ {m} + \mathrm{Y} _ {1} ^ {r _ {m}}) ^ {p _ {m}} - x _ {1} = 0.\tag{3}
\]

Let us write \(f(\mathbf{p}) = p_{1} + r_{2}p_{2} + \cdots + r_{m}p_{m}\) and suppose that the \(r_{i}\) are chosen so that the \(f(\mathbf{p})\) are distinct (it suffices for example to take \(r_{i} = h^{i}\) , where h is an integer strictly greater than all the \(p_{j}\) (Set Theory, Chapter III, §5, no. 7, Proposition 8)). Then there will be a unique system \(\mathbf{p} = (p_{1}, \ldots, p_{m})\) such that \(f(\mathbf{p})\) is maximal and relation (3) may be written

\[
a _ {\mathbf {p}} \mathrm{Y} _ {1} ^ {f (\mathbf {p})} + \sum_ {j <   f (\mathbf {p})} \mathrm{Q} _ {j} (x _ {1}, \dots , x _ {m}) \mathrm{Y} _ {1} ^ {j} = 0\tag{4}
\]

where the \(Q_{j}\) are polynomials in \(k[Y_{1},\ldots,Y_{m}]\) ; as \(a_{p} \neq 0\) is invertible in k, (4) is certainly an equation of integral dependence with coefficients in B, whence our assertion.

The field of fractions \(k(\mathrm{Y}_1, \ldots, \mathrm{Y}_m)\) of \(\mathbf{A}\) is therefore algebraic over the field of fractions \(k(x_1, \ldots, x_m)\) of \(\mathbf{B}\) , which proves (Algebra, Chapter V, § 5, no. 3, Theorem 4) that the \(x_i\) ( \(1 \leqslant i \leqslant m\) ) are algebraically independent. Moreover, \(a_1 \cap B = Bx_1\) ; for every element \(z \in a_1 \cap B\) may be written \(z = x_1z'\) where \(z' \in A \cap k(x_1, \ldots, x_m)\) ; but \(A \cap k(x_1, \ldots, x_m) = k[x_1, \ldots, x_m] = B\) since B is integrally closed (§ 1, no. 3, Corollary 2 to Proposition 13); therefore \(z' \in B\) , which completes the proof of properties (a) and (b) in this case.

\textbf{(A2) General case \((p = 1)\) .}\par

We argue by induction on \(m\) , the case \(m = 0\) being trivial. We may obviously suppose that \(a_1 \neq 0\) (otherwise we may take \(x_i = Y_i\) for \(1 \leqslant i \leqslant m\) and \(h(1) = 0\) ). Let \(x_1\) be a non-zero element of \(a_1\) ; by (A1) there exist \(t_2, \ldots, t_m\) such that \(x_1, t_2, \ldots, t_m\) are algebraically independent over \(k\) , A is integral over \(C = k[x_1, t_2, \ldots, t_m]\) and \(x_1A \cap C = x_1C\) . By the induction hypothesis there exist elements \(x_2, \ldots, x_m\) of \(k[t_2, \ldots, t_m]\) and an integer \(h\) such that \(k[t_2, \ldots, t_m]\) is integral over \(B' = k[x_2, \ldots, x_m]\) , \(x_2, \ldots, x_m\) are algebraically independent over \(k\) and the ideal \(a_1 \cap B'\) is generated by \(x_2, \ldots, x_h\) . Then C is integral over \(B = k[x_1, x_2, \ldots, x_m]\) ( \(\S\) 1, no. 1, Corollary to Proposition 5) and hence so is A ( \(\S\) 1, no. 1, Proposition 6); the same argument as in the case (A1) shows that \(x_1, \ldots, x_m\) are algebraically independent over \(k\) ; finally, as \(x_1 \in a_1\) and \(B = B'[x_1]\) , \(a_1 \cap B = Bx_1 + (a_1 \cap B')\) and, as \(a_1 \cap B'\) is generated (in \(B'\) ) by \(x_2, \ldots, x_h\) , \(a_1 \cap B\) is generated (in B) by \(x_1, x_2, \ldots, x_h\) .

\textbf{(B) Passage from p - 1 to p.}\par

Let \(t_1, \ldots, t_m\) be elements of A satisfying the conditions of the theorem for the increasing sequence of ideals \(a_1 \subset \ldots \subset a_{p-1}\) and let us write \(r = h(p - 1)\) . By (A2) there exist elements \(x_{r+1}, \ldots, x_m\) of \(k[t_{r+1}, \ldots, t_m]\) and an integer \(s\) such that

\[
\mathbf {C} = k [ t _ {r + 1}, \dots , t _ {m} ]
\]

is integral over \(\mathbf{B}' = k[x_{r+1},\ldots,x_m]\) , \(x_{r+1},\ldots,x_m\) are algebraically independent over \(k\) and the ideal \(\mathfrak{a}_p\cap\mathbf{B}'\) is generated by \(x_{r+1},\ldots,x_s\) . Writing \(x_i=t_i\) for \(i\leqslant r\) the family \((x_i)_{1\leqslant i\leqslant m}\) obtained solves the problem with \(h(p)=s\) . For A is integral over \(\mathbf{C}[t_1,\ldots,t_r]=\mathbf{C}[x_1,\ldots,x_r]\) and hence also over \(\mathbf{B}=k[x_1,\ldots,x_m]=\mathbf{B}'[x_1,\ldots,x_r]\) since C is integral over \(\mathbf{B}'\) (§ 1, no. 1, Corollary to Proposition 5 and Proposition 6); it can be shown as in the case (A1) that the \(x_i\) are algebraically independent over \(k\) . On the other hand, for \(j\leqslant p-1\) , the ideal

\[
a _ {j} \cap k [ x _ {1}, \dots , x _ {r}, t _ {r + 1}, \dots , t _ {m} ]
\]

is by hypothesis the set of polynomials in \(x_1, \ldots, x_r, t_{r+1}, \ldots, t_m\) all of whose monomials contain one of the elements \(x_1, \ldots, x_{h(j)}\) ; as \(x_{r+1}, \ldots, x_m\) are polynomials in \(t_{r+1}, \ldots, t_m\) with coefficients in \(k\) , it is seen immediately that a polynomial in \(x_1, \ldots, x_r, x_{r+1}, \ldots, x_m\) (with coefficients in \(k\) ) can belong to \(a_j\) only if all its monomials contain one of the elements \(x_1, \ldots, x_{h(j)}\) . Finally, as \(x_1, \ldots, x_r\) belong to \(a_{p-1}\) and hence also to \(a_p, a_p \cap B'[x_1, \ldots, x_r]\) consists of the polynomials in \(x_1, \ldots, x_r\) with coefficients in \(B'\) whose constant term belongs to \(a_p \cap B'\) ; this ideal is therefore generated by \(x_1, \ldots, x_r, x_{r+1}, \ldots, x_t\) .

COROLLARY 1. Let A be an integral domain and B a finitely generated A-algebra containing A as a subring. Then there exist an element \(s \neq 0\) of A and a subalgebra B' of B isomorphic to a polynomial algebra \(\mathrm{A}[\mathrm{Y}_{1}, \ldots, \mathrm{Y}_{n}]\) such that \(\mathrm{B}[s^{-1}]\) (Chapter II, §2, no. 1) is integral over \(\mathrm{B}'[s^{-1}]\) .

We write \(S = A - \{0\}\) and let \(k = S^{-1}A\) the field of fractions of \(A\) ; clearly \(S^{-1}B\) is a finitely generated \(k\) -algebra and, as it contains \(k\) by hypothesis (Chapter II, § 2, no. 4, Theorem 1), it is not reduced to 0. By Theorem 1 (applied to \(p = 1\) and \(a_1 = 0\) ) there exists therefore a finite sequence \((x_i)_{1 \leqslant i \leqslant n}\) of elements of \(S^{-1}B\) which are algebraically independent over \(k\) and such that \(S^{-1}B\) is integral over \(k[x_1, \ldots, x_n]\) . Let \((z_j)_{1 \leqslant j \leqslant m}\) be a system of generators of the A-algebra \(B\) ; in \(S^{-1}B\) each of the \(z_j/1\) satisfies an equation of integral dependence

\[
(z _ {j} / 1) ^ {q _ {j}} + \sum_ {h <   q _ {j}} \mathrm{P} _ {h j} (x _ {1}, \dots , x _ {n}) (z _ {j} / 1) ^ {h} = 0\tag{5}
\]

where the \(P_{hj}\) are polynomials in the \(x_{i}\) with coefficients in k. There exists an element \(s \neq 0\) of A such that we may write \(x_{i} = y_{i}/s\) where \(y_{i} \in B\) for \(1 \leqslant i \leqslant n\) and all the coefficients of the \(P_{hj}\) are of the form c/s where \(c \in A\) ; finally, replacing if need be, s by a product of elements of S, we may assume that in B

\[
s z _ {j} ^ {q _ {j}} + \sum_ {h <   q _ {j}} Q _ {h j} z _ {j} ^ {h} = 0\tag{6}
\]

where the \(Q_{hj}\) are polynomials in \(y_{1},\ldots,y_{n}\) with coefficients in A; if we write \(z_{j}^{\prime}=sz_{j}\) it is seen, multiplying (6) by \(s^{q_{j}-1}\) , that \(z_{j}^{\prime}\) is integral over \(B^{\prime}=A[y_{1},\ldots,y_{n}]\) . We show that the \(y_{i}\) are algebraically independent over A; if there is a relation of the form \(\sum_{p}a_{p}y_{1}^{p_{1}}\ldots y_{n}^{p_{n}}=0\) where \(a_{p}\in A\) for all p, we deduce that \(\sum_{p}a_{p}^{\prime}x_{1}^{p_{1}}\ldots x_{n}^{p_{n}}=0\) in \(S^{-1}B\) , where \(a_{p}^{\prime}=a_{p}s^{p_{1}+\cdots+p_{n}}\) in k; by hypothesis therefore \(a_{p}^{\prime}=0\) for all p, whence \(a_{p}=0\) for all p.~Moreover, in the ring \(B[s^{-1}]\) each of the \(z_{j}^{\prime}/1\) is integral over \(B^{\prime}[s^{-1}]\) (§1, no. 1, Proposition 2) and, as \(z_{j}/1=(z_{j}^{\prime}/1)(1/s)\) in \(B[s^{-1}]\) , it is seen that the \(z_{j}/1\) are integral over \(B^{\prime}[s^{-1}]\) , which completes the proof (§1, no. 1, Proposition 4).

COROLLARY 2. Let K be a field, A a subring of K and L the field of fractions of A. If K is a finitely generated A-algebra, {[}K:L{]} is finite and there exists \(a \neq 0\) in A such that \(\mathbf{L} = \mathbf{A}[a^{-1}]\) .

It follows from Corollary 1 that there exist elements \(x_{1}, \ldots, x_{n}\) of \(\mathbf{K}\) and an element \(a \neq 0\) of \(\mathbf{A}\) such that \(x_{1}, \ldots, x_{n}\) are algebraically independent over A (and therefore over L) and that K is integral over the subring \(A[x_{1},\ldots,x_{n},a^{-1}]\) . Then it follows from §2, no. 1, Lemma 2 that \(A[x_{1},\ldots,x_{n},a^{-1}]\) is a field. But the only invertible elements of a polynomial ring \(C[Y_{1},\ldots,Y_{n}]\) over an integral domain C are the invertible elements of C; applying this remark to \(C = A[a^{-1}]\) , it is seen that necessarily n = 0 and that \(A[a^{-1}]\) is a field equal to L by definition of the latter. As K is integral over L and is a finitely generated L-algebra, the degree {[}K:L{]} is finite (§1, no. 1, Proposition 4).

COROLLARY 3. Let A be an integral domain, B a finitely generated A-algebra and b an element of B such that \(zb^n \neq 0\) for all \(z \neq 0\) in A and every integer \(n > 0\) . Let \(\rho: A \to B\) be the canonical homomorphism; there exists \(a \neq 0\) in A such that, for every homomorphism \(f\) from A to an algebraically closed field L such that \(f(a) \neq 0\) , there exists a homomorphism \(g\) from B to L such that \(g(b) \neq 0\) and \(f = g \circ \rho\) .

The hypothesis on \(b\) implies that, if \(h\) is the canonical homomorphism \(x \mapsto x / 1\) of \(\mathbf{B}\) to \(\mathbf{B}[b^{-1}]\) , the homomorphism \(h \circ \rho\) of \(\mathbf{A}\) to \(\mathbf{B}[b^{-1}]\) is injective. By Corollary 1 there therefore exist an element \(a \neq 0\) of \(\mathbf{A}\) and a subring \(\mathbf{B}'\) of \(\mathbf{B}[b^{-1}]\) such that \(\mathbf{B}[b^{-1}, a^{-1}]\) is integral over \(\mathbf{B}'[a^{-1}]\) and \(\mathbf{B}'\) is isomorphic to a polynomial algebra \(\mathbf{A}[\mathbf{Y}_1, \ldots, \mathbf{Y}_n]\) . Let \(f\) be a homomorphism from \(\mathbf{A}\) to an algebraically closed field \(\mathbf{L}\) such that \(f(a) \neq 0\) ; there exists a homomorphism from \(\mathbf{A}[\mathbf{Y}_1, \ldots, \mathbf{Y}_n]\) to \(\mathbf{L}\) extending \(f\) and hence there exists a homomorphism \(f'\) from \(\mathbf{B}'\) to \(\mathbf{L}\) extending \(f\) . As \(f'(a) \neq 0\) in \(\mathbf{L}\) , there exists a homomorphism \(f''\) from \(\mathbf{B}'[a^{-1}]\) to \(\mathbf{L}\) such that

\[
f ^ {\prime \prime} (x / a ^ {n}) = f ^ {\prime} (x). (f (a)) ^ {- n}
\]

for all \(x \in \mathbf{B}'\) and all \(n > 0\) (Chapter II, §2, no. 1, Proposition 1). Finally, as \(\mathbf{B}[b^{-1}, a^{-1}]\) is integral over \(\mathbf{B}'[a^{-1}]\) , there exists a homomorphism \(f'''\) from \(\mathbf{B}[b^{-1}, a^{-1}]\) to \(\mathbf{L}\) extending \(f''\) (§2, no. 1, Corollary 4 to Theorem 1). If \(j: x \mapsto x / 1\) is the canonical homomorphism from \(\mathbf{B}\) to \(\mathbf{B}[b^{-1}, a^{-1}]\) , \(g = f''' \circ j\) solves the problem for \(j(b)\) is invertible in \(\mathbf{B}[b^{-1}, a^{-1}]\) and hence \(f'''(j(b)) \neq 0\) in \(\mathbf{L}\) .

Note that, if B is assumed to be an integral domain and A ⊆ B in Corollary 3, the hypothesis on b is equivalent to \(b \neq 0\) .

\subsection{The integral closure of a finitely generated algebra over a field}

THEOREM 2. Let A be a finitely generated integral k-algebra, K its field of fractions and A' the integral closure of A in a field K' which is a finite algebraic extension of K. Then A' is a finitely generated A-module and a finitely generated k-algebra.

By Theorem 1 there exists a subalgebra C of A isomorphic to a polynomial algebra \(k[\mathbf{X}_1, \ldots, \mathbf{X}_n]\) and such that A is integral over C; \(\mathbf{A}'\) is obviously the integral closure of C in \(\mathbf{K}'\) (§ 1, no. 1, Proposition 6); we may therefore confine our attention to the case where \(A = k[X_{1}, \ldots, X_{n}]\) . Let N be the quasi-Galois extension of \(K'\) (in an algebraic closure of \(K'\) ) generated by \(K'\) , which is a finite algebraic extension of K (Algebra, Chapter V, § 6, no. 3, Corollary 1 to Proposition 9). It will suffice to prove that the integral closure B of A in N is a finitely generated A-module, for \(A'\) is a sub-A-module of B and A is a Noetherian ring (Chapter III, § 2, no. 10, Corollary 2 to Theorem 2). We may therefore confine our attention to the case where \(K'\) is a quasi-Galois extension of K. Then we know (Algebra, Chapter V, § 10, no. 9, Proposition 14) that \(K'\) is a (finite) Galois extension of a (finite) radicial extension \(K''\) of K. If \(A''\) is the integral closure of A in \(K''\) , \(A'\) is the integral closure of \(A''\) in \(K'\) and it will suffice to prove that \(A''\) is a finitely generated A-module and \(A'\) is a finitely generated \(A''\) -module. Now, if it has been proved that \(A''\) is a finitely generated A-module, it is a Noetherian domain, integrally closed by definition; the fact that \(A'\) is a finitely generated \(A''\) -module will follow from § 1, no. 6, Corollary 1 to Proposition 18.

We see therefore that we may confine our attention to the case where \(\mathbf{A} = k[\mathbf{X}_1, \ldots, \mathbf{X}_n]\) and where \(\mathbf{K}'\) is a finite radicial extension of \(\mathbf{K} = k(\mathbf{X}_1, \ldots, \mathbf{X}_n)\) . Then \(\mathbf{K}'\) is generated by a finite family of elements \((y_i)_{1 \leqslant i \leqslant m}\) and there exists a power \(q\) of the characteristic exponent of \(k\) such that \(y_i^q \in k(\mathbf{X}_1, \ldots, \mathbf{X}_n)\) . Let \(c_j\) ( \(1 \leqslant j \leqslant r\) ) be the coefficients of the numerators and denominators of the rational functions in \(\mathbf{X}_1, \ldots, \mathbf{X}_n\) equal to \(y_i^q\) ( \(1 \leqslant i \leqslant m\) ). Then \(\mathbf{K}'\) is contained in the extension \(\mathbf{L} = k'(\mathbf{X}_1^{q^{-1}}, \ldots, \mathbf{X}_n^{q^{-1}})\) , where \(k' = k(c_1^{q^{-1}}, \ldots, c_n^{q^{-1}})\) (we are in an algebraic closure of \(\mathbf{K}'\) ) and \(\mathbf{A}'\) is contained in the algebraic closure \(\mathbf{B}'\) of \(\mathbf{A}\) in \(\mathbf{L}'\) . Now, \(k'\) is algebraic over \(k\) and hence \(\mathbf{C}' = k'[\mathbf{X}_1, \ldots, \mathbf{X}_n]\) is integral over \(\mathbf{A}\) ( \(\S 1\) , no. 1, Proposition 5); as \(k'[\mathbf{X}_1^{q^{-1}}, \ldots, \mathbf{X}_n^{q^{-1}}]\) is integrally closed ( \(\S 1\) , no. 3, Corollary 2 to Proposition 13), it is seen that this ring is the integral closure of \(\mathbf{C}'\) in \(\mathbf{L}'\) and hence also that of \(\mathbf{A}\) ( \(\S 1\) , no. 1, Proposition 6), in other words \(\mathbf{B}' = k'[\mathbf{X}_1^{q^{-1}}, \ldots, \mathbf{X}_n^{q^{-1}}]\) . Now clearly \(\mathbf{B}'\) is a finitely generated \(\mathbf{C}'\) -module ( \(\S 1\) , no. 1, Proposition 4) and, as \(k'\) is a finite extension of \(k\) , \(\mathbf{C}'\) is a finitely generated A-module and hence \(\mathbf{B}'\) is a finitely generated A-module; since \(\mathbf{A}\) is Noetherian and \(\mathbf{A}' \subset \mathbf{B}'\) , \(\mathbf{A}'\) is a finitely generated A-module.

\subsection{The Nullstellensatz}

PROPOSITION 1. Let A be a finitely generated algebra over a field \(k\) and \(L\) the algebraic closure of \(k\) .

\begin{enumerate}
\def\labelenumi{(\roman{enumi})}
\item
  If \(\mathbf{A} \neq \{0\}\) , there exists a \(k\) -homomorphism from \(\mathbf{A}\) to \(\mathbf{L}\) .
\item
  Let \(f_{1}, f_{2}\) be two \(k\) -homomorphisms from \(\mathbf{A}\) to \(\mathbf{L}\) . For \(f_{1}\) and \(f_{2}\) to have the same kernel, it is necessary and sufficient that there exist a \(k\) -automorphism \(s\) of \(\mathbf{L}\) such that \(f_{2} = s \circ f_{1}\) .
\item
  Let \(\mathfrak{a}\) be an ideal of \(\mathbf{A}\) . For \(\mathfrak{a}\) to be maximal, it is necessary and sufficient that it be the kernel of a \(k\) -homomorphism from \(\mathbf{A}\) to \(\mathbf{L}\) .
\item
  For an element \(x\) of \(\mathbf{A}\) to be such that \(f(x) = 0\) for every \(k\) -homomorphism \(f\) from \(\mathbf{A}\) to \(\mathbf{L}\) , it is necessary and sufficient that \(x\) be nilpotent.
\end{enumerate}

Assertion (i) follows from no. 1, Corollary 3 to Theorem 1 applied replacing A by k, B by A, b by the unit element of B and f by the canonical injection of k into L.

If \(f\) is a \(k\) -homomorphism from \(\mathbf{A}\) to \(\mathbf{L}\) , \(f(\mathbf{A})\) is a subring of \(\mathbf{L}\) containing \(k\) ; as \(\mathbf{L}\) is an algebraic extension of \(k\) , \(f(\mathbf{A})\) is a field (Algebra, Chapter V, § 3, no. 2, Proposition 3) and, if \(\mathfrak{a}\) is the kernel of \(f\) , \(\mathbf{A} / \mathfrak{a}\) , isomorphic to \(f(\mathbf{A})\) , is therefore a field, which proves that \(\mathfrak{a}\) is maximal. Conversely, if \(\mathfrak{a}\) is a maximal ideal of \(\mathbf{A}\) , it follows from (i) that there exists a \(k\) -homomorphism from \(\mathbf{A} / \mathfrak{a}\) to \(\mathbf{L}\) and hence a \(k\) -homomorphism of \(\mathbf{A}\) to \(\mathbf{L}\) whose kernel \(\mathfrak{b}\) contains \(\mathfrak{a}\) ; but as \(\mathfrak{a}\) is maximal, \(\mathfrak{b} = \mathfrak{a}\) ; this proves (iii).

We now prove (ii). If \(s\) is a \(k\) -automorphism of \(L\) such that \(f_2 = s \circ f_1\) , clearly \(f_1\) and \(f_2\) have the same kernel. Conversely, suppose that \(f_1\) and \(f_2\) have the same kernel; then there exists a \(k\) -isomorphism \(s_0\) of the field \(f_1(A)\) onto the field \(f_2(A)\) such that \(f_2 = s_0 \circ f_1\) ; but by Algebra, Chapter V, § 6, no. 3, Proposition 7, \(s_0\) extends to a \(k\) -automorphism \(s\) of \(L\) and hence \(f_2 = s \circ f_1\) .

Finally, if \(x \in \mathbf{A}\) is such that \(x^n = 0\) , for every \(k\) -homomorphism \(f\) from \(\mathbf{A}\) to \(\mathbf{L}\) , \((f(x))^n = f(x^n) = 0\) and hence \(f(x) = 0\) since \(\mathbf{L}\) is a field. Conversely, suppose that \(x \in \mathbf{A}\) is not nilpotent; then \(\mathbf{A}[x^{-1}]\) is a finitely generated A-algebra (and therefore a finitely generated \(k\) -algebra) not reduced to 0 (Chapter II, §2, no. 1, Remark 3) and hence there exists a \(k\) -homomorphism \(g\) from \(\mathbf{A}[x^{-1}]\) to \(\mathbf{L}\) by (i). If \(j: \mathbf{A} \to \mathbf{A}[x^{-1}]\) is the canonical homomorphism, \(f = g \circ j\) is a \(k\) -homomorphism from \(\mathbf{A}\) to \(\mathbf{L}\) and \(f(x)g(1/x) = g(x/1)g(1/x) = g(1) = 1\) , whence \(f(x) \neq 0\) .

Let \(k\) be a field and \(\mathbf{L}\) an extension field of \(k\) ; an element \(\mathbf{x} = (x_1, \ldots, x_n)\) of \(\mathbf{L}^n\) is called a zero in \(\mathbf{L}^n\) of an ideal \(\mathfrak{r}\) of the polynomial ring \(k[\mathbf{X}_1, \ldots, \mathbf{X}_n]\) if

\[
\mathrm{P} (\mathbf {x}) = \mathrm{P} (x _ {1}, \dots , x _ {n}) = 0
\]

for all \(\mathbf{P} \in \mathfrak{r}\) .

LEMMA 1. Let A be a finitely generated algebra over a field \(k\) , \((a_i)_{1 \leqslant i \leqslant n}\) a system of generators of this algebra and \(\mathfrak{r}\) the ideal of algebraic relations between the \(a_i\) with coefficients in \(k\) (Algebra, Chapter IV, §2, no. 1). For every extension field \(\mathbf{L}\) of \(k\) , the mapping \(f \mapsto (f(a_i))_{1 \leqslant i \leqslant n}\) is a bijection of the set of \(k\) -homomorphisms from A to L onto the set of zeros of \(\mathfrak{r}\) in \(L^n\) .

There exists a unique \(k\) -algebra homomorphism \(h\) of \(k[\mathbf{X}_1, \ldots, \mathbf{X}_n]\) onto \(\mathbf{A}\) such that \(h(\mathbf{X}_i) = a_i\) for \(1 \leqslant i \leqslant n\) and by definition \(\mathfrak{r}\) is the kernel of \(h\) . The mapping \(f \mapsto f \circ h\) is a bijection of the set of \(k\) -homomorphisms from \(A\) to \(L\) onto the set of \(k\) -homomorphisms from \(k[X_1, \ldots, X_n]\) to \(L\) which are zero on \(r\) . For every polynomial \(P \in k[X_1, \ldots, X_n]\) and every element \(\mathbf{x} = (x_1, \ldots, x_n) \in L^n\) we write \(h_{\mathbf{x}}(P) = P(\mathbf{x})\) ; then the mapping \(\mathbf{x} \mapsto h_{\mathbf{x}}\) is a bijection of \(L^n\) onto the set of \(k\) -homomorphisms from \(k[X_1, \ldots, X_n]\) to \(L\) (such a homomorphism being determined by its values at the \(X_i\) ( \(1 \leqslant i \leqslant n\) )); to say that \(h_{\mathbf{x}}\) is zero on \(r\) means that \(\mathbf{x}\) is a zero of \(r\) in \(L^n\) , whence the lemma.

If Proposition 1 is applied to the algebra \(\mathbf{A} = k[\mathbf{X}_1, \ldots, \mathbf{X}_n] / \mathfrak{r}\) , where \(\mathfrak{r}\) is an ideal of \(k[\mathbf{X}_1, \ldots, \mathbf{X}_n]\) distinct from the whole ring, we obtain by Lemma 1 the following statement:

PROPOSITION 2 (Hilbert's Nullstellensatz). Let \(k\) be a field and \(L\) an algebraic closure of \(k\) .

\begin{enumerate}
\def\labelenumi{(\roman{enumi})}
\item
  Every ideal \(\mathfrak{r}\) of \(k[\mathbf{X}_1, \ldots, \mathbf{X}_n]\) not containing 1 admits at least one zero in \(\mathbf{L}^n\) .
\item
  Let \(\mathbf{x} = (x_1, \ldots, x_n)\) , \(\mathbf{y} = (y_1, \ldots, y_n)\) be two elements of \(\mathbf{L}^n\) ; for the set of polynomials of \(k[\mathbf{X}_1, \ldots, \mathbf{X}_n]\) zero at \(\mathbf{x}\) to be identical with the set of polynomials of \(k[\mathbf{X}_1, \ldots, \mathbf{X}_n]\) zero at \(\mathbf{y}\) , it is necessary and sufficient that there exists a \(k\) -automorphism \(s\) of \(\mathbf{L}\) such that \(y_i = s(x_i)\) for \(1 \leqslant i \leqslant n\) .
\item
  For an ideal \(\mathfrak{a}\) of \(k[\mathbf{X}_1, \ldots, \mathbf{X}_n]\) to be maximal, it is necessary and sufficient that there exist an \(\mathbf{x}\) in \(\mathbf{L}^n\) such that \(\mathfrak{a}\) is the set of polynomials of \(k[\mathbf{X}_1, \ldots, \mathbf{X}_n]\) zero at \(\mathbf{x}\) .
\item
  For a polynomial \(\mathbf{Q}\) of \(k[\mathbf{X}_1, \ldots, \mathbf{X}_n]\) to be zero on the set of zeros in \(\mathbf{L}^n of\) an ideal \(\mathfrak{r}\) of \(k[\mathbf{X}_1, \ldots, \mathbf{X}_n]\) , it is necessary and sufficient that there exist an integer \(m > 0\) such that \(\mathbf{Q}^m \in \mathfrak{r}\) .
\end{enumerate}

\subsection{Jacobson rings}

DEFINITION 1. A ring \(\mathbf{A}\) is called a Jacobson ring if every prime ideal of \(\mathbf{A}\) is the intersection of a family of maximal ideals.

\textbf{Examples}\par

\begin{enumerate}
\def\labelenumi{(\arabic{enumi})}
\item
  Every field is a Jacobson ring.
\item
  The ring Z is a Jacobson ring, the unique prime ideal which is not maximal (0) being the intersection of the maximal ideals \((p)\) of Z, where p runs through the set of prime numbers (cf.~Proposition 4).
\item
  Let A be a Jacobson ring and let a be an ideal of A. Then A/a is a Jacobson ring, for the ideals of A/a are of the form b/a, where b is an ideal of A containing a and b/a is prime (resp. maximal) if and only if b is.
\end{enumerate}

PROPOSITION 3. For a ring A to be a Jacobson ring, it is necessary and sufficient that, for every ideal a of A, the Jacobson radical of A/a be equal to its nilradical (Chapter II, § 2, no. 6).

The Jacobson radical (resp. nilradical) of A/a is the intersection of the maximal (resp. prime) ideals of A/a (Algebra, Chapter VIII, § 6, no. 3, Definition 3 and Commutative Algebra, Chapter II, § 2, no. 6, Proposition 13). The stated condition means therefore that for every ideal a of A the intersection of the prime ideals containing a is equal to the intersection of the maximal ideals containing a. This condition obviously holds for every ideal a of A if A is a Jacobson ring; conversely, if it holds for every prime ideal of A, A is a Jacobson ring by definition.

COROLLARY. Let A be a Jacobson ring; for every ideal a of A, the radical of a is the intersection of the maximal ideals of A containing a.

It is sufficient to note that A/a is a Jacobson ring.

PROPOSITION 4. Let A be a principal ideal domain and \((p_{\lambda})_{\lambda \in L}\) a representative system of extremal elements of A (Algebra, Chapter VII, § 1, no. 3, Definition 2). For A to be a Jacobson ring, it is necessary and sufficient that L be infinite.

The maximal ideals of A are the \(Ap_{\lambda}\) (loc. cit., no. 2, Proposition 2). If L is finite, their intersection is the ideal Ax, where \(x = \prod_{\lambda \in L} p_{\lambda}\) (ibid.) and hence different from (0); on the other hand, if L is infinite, the intersection of the \(Ap_{\lambda}\) is (0), every element \(\neq 0\) of A being divisible by only a finite number of extremal elements (loc. cit., no. 3, Theorem 2). The proposition then follows from the fact that (0) is the only prime ideal which is not maximal in A (Algebra, Chapter VI, §1, no. 13, Proposition 14 (DIV)).

PROPOSITION 5. Let A be a ring and B an A-algebra integral over A. If A is a Jacobson ring, so is B.

Replacing A by its canonical image in B, we may assume that \(A \subset B\) . Let \(p'\) be a prime ideal of B and let \(p = A \cap p'\) . There exists by hypothesis a family \((m_{\lambda})_{\lambda \in L}\) of maximal ideals of A whose intersection is equal to p.~For all \(\lambda \in L\) there exists a maximal ideal \(m_{\lambda}'\) of B lying above \(m_{\lambda}\) and containing \(p'\) (§ 2, no. 1, Proposition 1 and Corollary 2 to Theorem 1). If we write \(q' = \bigcap_{\lambda \in L} m_{\lambda}'\) , then \(q' \cap A = \bigcap_{\lambda \in L} m_{\lambda} = p\) and \(q' \supset p'\) , whence \(q' = p'\) (§ 2, no. 1, Corollary 1 to Proposition 1).

THEOREM 3. Let A be a Jacobson ring, B a finitely generated A-algebra and \(\rho: \mathbf{A} \to \mathbf{B}\) the canonical homomorphism. Then:

\begin{enumerate}
\def\labelenumi{(\roman{enumi})}
\item
  B is a Jacobson ring.
\item
  For every maximal ideal \(\mathfrak{m}'\) of \(\mathbf{B}\) , \(\mathfrak{m} = \rho^{-1} (\mathfrak{m}')\) is a maximal ideal of \(\mathbf{A}\) and \(\mathbf{B} / \mathfrak{m}'\) is a finite algebraic extension of \(\mathbf{A} / \mathfrak{m}\) .
\end{enumerate}

Let \(\mathfrak{p}'\) be a prime ideal of B and \(\mathfrak{p} = \rho^{-1} (\mathfrak{p}')\) . Let \(v\) be an element \(\neq 0\) of \(\mathrm{B / p'}\) .

As \(B/p'\) is a finitely generated integral \((A/p)\) -algebra and the canonical homomorphism \(\phi: A/p \to B/p'\) is injective, there exists an element \(u \neq 0\) of \(A/p\) such that, for every homomorphism \(f\) from \(A/p\) to an algebraically closed field \(L\) whose kernel does not contain \(u\) , there exists a homomorphism \(g\) from \(B/p'\) to \(L\) whose kernel does not contain \(v\) and for which \(f = g \circ \phi\) (no. 1, Corollary 3 to Theorem 1). Since \(A\) is a Jacobson ring, there exists a maximal ideal \(m\) of \(A\) containing \(p\) and such that \(u \notin m/p\) . We take \(L\) to be an algebraic closure of \(A/m\) and \(f\) to be the canonical homomorphism \(A/p \to L\) ; let

\[
g \colon \mathrm{B} / \mathfrak {p} ^ {\prime} \rightarrow \mathrm{L}
\]

be a homomorphism such that \(f = g \circ \phi\) and \(g(v) \neq 0\) . Then

\[
\mathbf {A} / \mathfrak {m} \subset g (\mathbf {B} / \mathfrak {p} ^ {\prime}) \subset \mathbf {L},
\]

hence \(g(\mathbf{B} / \mathfrak{p}')\) is a subfield of \(\mathbf{K}\) (Algebra, Chapter V, § 3, no. 2, Proposition 3) and the kernel of \(g\) is therefore a maximal ideal of \(\mathbf{B} / \mathfrak{p}'\) not containing \(v\) . Thus it is seen that the intersection of the maximal ideals of \(\mathbf{B} / \mathfrak{p}'\) is reduced to 0, which proves that \(\mathbf{B}\) is a Jacobson ring. Moreover, if \(\mathfrak{p}'\) is maximal, \(g\) is necessarily injective and hence \(\mathfrak{p} = \mathfrak{m}\) is maximal; finally \(\mathbf{B} / \mathfrak{p}'\) is then a finitely generated algebra over the field \(\mathbf{A} / \mathfrak{m}\) and hence is a finite extension of \(\mathbf{A} / \mathfrak{m}\) (no. 1, Corollary 2 to Theorem 1).

COROLLARY 1. Every finitely generated algebra A over Z is a Jacobson ring; for a prime ideal p of A to be maximal, it is necessary and sufficient that the ring A/p be finite.

If the integral domain A/p is finite, it is a field, as, for every \(u \neq 0\) in A/p, the mapping \(v \mapsto uv\) of A/p to itself is injective and hence bijective since A/p is finite. Conversely, for every maximal ideal m of A, the inverse image of m in Z is a maximal ideal (p) and A/m is finite over the prime field \(\mathbf{Z}/(\mathfrak{p}) = \mathbf{F}_{\mathfrak{p}}\) by Theorem 3.

COROLLARY 2. Let \((\mathbf{P}_{\lambda})_{\lambda \in \mathbf{L}}\) be a family of polynomials in \(\mathbf{Z}[\mathbf{X}_1,\ldots ,\mathbf{X}_n]\) and let \(\mathbf{Q}\) be a polynomial in \(\mathbf{Z}[\mathbf{X}_1,\dots ,\mathbf{X}_n]\) such that, for every system of elements \((x_i)_{1\leqslant i\leqslant n}\) belonging to a finite field and for which \(\mathrm{P}_{\lambda}(x_1,\ldots ,x_n) = 0\) for all \(\lambda\) , also \(\mathrm{Q}(x_1,\ldots ,x_n) = 0\) . Then, if \(\mathfrak{a}\) is the ideal of \(\mathbf{Z}[\mathbf{X}_1,\ldots ,\mathbf{X}_n]\) generated by the \(\mathrm{P}_{\lambda}\) , there exists an integer \(m > 0\) such that \(\mathbf{Q}^m\in \mathfrak{a}\) . Moreover, for every reduced ring R and every system \((y_i)_{1\leqslant i\leqslant n}\) of elements of R such that \(\mathrm{P}_{\lambda}(y_1,\ldots ,y_n) = 0\) for all \(\lambda\) , also \(\mathrm{Q}(y_1,\ldots ,y_n) = 0\) .

The second assertion follows from the first since the ideal of \(\mathbf{Z}[\mathrm{X}_1,\ldots ,\mathrm{X}_n]\) consisting of the polynomials \(\mathbf{P}\) such that \(\mathrm{P}(y_1,\dots,y_n) = 0\) contains a. To show the first assertion, it suffices to note that, for every maximal ideal m of \(\mathbf{A} = \mathbf{Z}[\mathrm{X}_1,\dots,\mathrm{X}_n]\) containing a, A/m is a finite field (Corollary 1) and the hypothesis implies that the canonical image of Q in A/m is zero; then Q belongs to the intersection of the maximal ideals of A containing a, which is the radical of a (Corollary to Proposition 3).

COROLLARY 3. Let A be a Jacobson ring. If there exists a finitely generated A-algebra B containing A and which is a field, then A is a field and B is an algebraic extension of A.

It suffices to apply Theorem 3 (ii) with \(\mathfrak{m}' = (0)\) .


\exercisesection{1}

\begin{enumerate}
\def\labelenumi{\arabic{enumi}.}
\item
  Let \(d\) be a rational integer not divisible by a square in \(\mathbf{Z}\) . Show that the elements of the field \(\mathbf{K} = \mathbf{Q}(\sqrt{d})\) which are integral over \(\mathbf{Z}\) are of the form \(a + b\sqrt{d}\) where \(a \in \mathbf{Z}\) , \(b \in \mathbf{Z}\) if \(d - 1 \not\equiv 0 \pmod{4}\) and the elements \((a + b\sqrt{d})/2\) where \(a \in \mathbf{Z}\) , \(b \in \mathbf{Z}\) , \(a\) and \(b\) both even or \(a\) and \(b\) both odd, if \(d - 1 \equiv 0 \pmod{4}\) (cf.~Algebra, Chapter VII, § 1, Exercise 8). Deduce that \(\mathbf{A} = \mathbf{Z}[\sqrt{5}]\) is not integrally closed and give an example of an element of \(\mathbf{Q}[\sqrt{5}]\) which is integral over \(\mathbf{A}\) but whose minimal polynomial over the field of fractions of \(\mathbf{A}\) does not have coefficients in \(\mathbf{A}\) .
\item
  Let K be a commutative field and A the sub-K-algebra of the polynomial algebra \(K[X, Y]\) generated by the monomials \(Y^{k}X^{k+1} (k \geqslant 0)\) . Show that XY is such that A{[}XY{]} is contained in a finitely generated A-module but XY is not integral over A.
\item
  Give an example of an infinite sequence \((\mathbf{K}_{n})\) of extensions of a commutative field K of finite degree over K such that the K-algebra \(\prod_{n} K_{n}\) is not algebraic over K.
\item
  In the matrix ring \(\mathbf{M}_{2}(\mathbf{Q})\) give an example of two elements integral over Z but neither the sum nor the product of which is integral over Z (consider matrices of the form \(I + N\) , where N is nilpotent).
\item
  Let A be a commutative ring, B a commutative ring containing A and x an invertible element of B. Show that every element of \(A[x] \cap A[x^{-1}]\) is integral over A. (If \(y = a_{0} + \cdots + a_{p}x^{-p} = b_{0} + \cdots + b_{q}x^{q}\) , where the \(a_{i}\) and \(b_{j}\) are in A, show that the sub-A-module of B generated by 1, \(x, \ldots, x^{p+q+1}\) is a faithful A{[}y{]}-module.)
\item
  Let A be a commutative ring and B a commutative A-algebra; suppose that the minimal prime ideals of B are finite in number; let \(p_{i}\) ( \(1 \leqslant i \leqslant n\) ) be these ideals (note that this hypothesis holds if B is Noetherian). For an element \(x \in B\) to be integral over A, it is necessary and sufficient that each of its canonical images \(x_{i}\) in \(B/p_{i}\) be integral over A. (If this condition holds, show first that x is integral over the subalgebra of B generated by the nilradical \(N = \bigcap_{i} p_{i}\) of B and use the fact that the elements of N are nilpotent.)
\item
  Give an example of a reduced ring A which is integrally closed in its total ring of fractions and which is not a product of integral domains (take the integral closure of a field K in the infinite product \(\prod_{n} K_{n}\) of suitable algebraic extensions of K).
\item
  Let A be a commutative ring and B a finitely generated A-algebra. Show that there exists a finite A-algebra C which is a free A-module and a surjective A-homomorphism \(C \rightarrow B\) .
\item
  Show that the quotient domain \(\mathbf{Z}[\mathbf{X}] / (\mathbf{X}^2 +4)\) is not integrally closed even though \(\mathbf{Z}[\mathbf{X}]\) is.
\item
  Let A be a completely integrally closed domain. Show that for every set I the polynomial domain \(A[X_{i}]_{i \in I}\) and the domain of formal power series \(A[[X_{i}]]_{i \in I}\) (Algebra, Chapter IV, §5, Exercise 1) are completely integrally closed. (Use Proposition 4 of no. 4 and the following lemma: for every finite subset J of I and all \(P \in A[[X_{i}]]_{i \in J}\) , let \(P_{J}\) be the formal power series in \(A[[X_{i}]]_{i \in J}\) derived from P by replacing by 0 in P all the \(X_{i}\) of index \(i \notin J\) ; if P, Q are two formal power series such that \(P_{J}\) is divisible by \(Q_{J}\) in \(A[[X_{i}]]_{i \in J}\) for all J, then P is divisible by Q in \(A[[X_{i}]]_{i \in I}\) . For this note that, if J, \(J'\) are two finite subsets of I such that \(J \subset J'\) and \(P_{J'} = L'Q_{J'}\) and \(P_{J} = LQ_{J}\) , then L is obtained from \(L'\) by replacing by 0 in \(L'\) the \(X_{i}\) of index \(i \in J' - J\) .)
\item
  \begin{enumerate}
  \def\labelenumii{(\alph{enumii})}
  \tightlist
  \item
    Let R be an integral domain, \((\mathbf{A}_{\alpha})\) a right directed family of subrings of R and A the union of the \(A_{\alpha}\) . Show that, if each of the \(A_{\alpha}\) is integrally closed, so is A.
  \end{enumerate}
\end{enumerate}

\begin{enumerate}
\def\labelenumi{(\alph{enumi})}
\setcounter{enumi}{1}
\tightlist
\item
  Let K be a commutative field and \(\mathbf{R} = \mathbf{K}(\mathbf{X}, \mathbf{Y})\) the field of rational functions in two indeterminates over K. For every integer n \textgreater{} 0 let \(A_{n}\) denote the subring \(\mathbf{K}[\mathbf{X}, \mathbf{Y}, \mathbf{Y}\mathbf{X}^{-n}]\) of R. Show that the \(A_{n}\) are completely integrally closed but that their union A is not. (For an element of the form \(\mathbf{P}(\mathbf{X}, \mathbf{Y}) \cdot \mathbf{X}^{-ns}\) , where \(P \in K[X, Y]\) , to belong to \(A_{n}\) show that it is necessary and sufficient that P belong to \(S^{s}\) , where S denotes the ideal of \(K[X, Y]\) generated by \(X^{n}\) and Y.)
\end{enumerate}

¶ * 12. Let A be the ring of integral functions (with complex values) of a complex variable. It is an integral domain whose field of fractions is the field of meromorphic functions of a complex variable.

\begin{enumerate}
\def\labelenumi{(\alph{enumi})}
\item
  Show that \(\mathbf{A}\) is completely integrally closed.
\item
  For all \(x \in \mathbf{A}\) , let \(Z(x)\) be the set of zeros of \(x\) (which is a closed subset of \(\mathbf{C}\) , which is discrete (and therefore countable) if \(x \neq 0\) ). If \(\mathfrak{a}\) is an ideal of \(\mathbf{A}\) distinct from \(\mathbf{A}\) and (0), show that the \(Z(x)\) for \(x \in \mathfrak{a}\) form a filter base \(\mathfrak{F}_{\mathfrak{a}}\) . (If \(x\) and \(y\) are two integral functions with no zero in common, show, using the Mittag-Leffler Theorem, that it is possible to write \(1 / xy = u / x + v / y\) , where \(u\) and \(v\) are in \(\mathbf{A}\) , and deduce that \(Ax + Ay = A\) ). Conversely, for every filter \(\mathfrak{F}\) on \(\mathbf{C}\) to which a closed discrete subset of \(\mathbf{C}\) belongs, the set \(\mathfrak{a}(\mathfrak{F})\) of \(x \in \mathbf{A}\) for which \(Z(x) \in \mathfrak{F}\) is an ideal of \(\mathbf{A}\) . Show that \(\mathfrak{F}_{\mathfrak{a}(\mathfrak{F})} = \mathfrak{F}\) for every filter \(\mathfrak{F}\) to which a closed discrete subset of \(\mathbf{C}\) belongs and \(\mathfrak{a}(\mathfrak{F}) \supset \mathfrak{a}\) for every ideal \(\mathfrak{a}\) of \(\mathbf{A}\) distinct from \(\mathbf{A}\) and (0).
\item
  If \(\mathfrak{p} \neq (0)\) is a prime ideal of \(\mathbf{A}\) , show that \(\mathfrak{F}_{\mathfrak{p}}\) is an ultrafilter. (Restricting it to the case where all the sets of \(\mathfrak{F}_{\mathfrak{p}}\) are infinite, show that, if \(\mathbf{M} \in \mathfrak{F}_{\mathfrak{p}}\) is the union of two infinite sets \(\mathbf{M}'\) , \(\mathbf{M}''\) with no point in common, one of these sets belongs to \(\mathfrak{F}_{\mathfrak{p}}\) .) Conversely, for every ultrafilter \(\mathfrak{U}\) on \(\mathbf{C}\) to which a closed discrete subset of \(\mathbf{C}\) belongs, \(\mathfrak{a}(\mathfrak{U})\) is a maximal ideal of \(\mathbf{A}\) ; obtain the converse.
\item
  Let \(\mathfrak{U}\) be an ultrafilter on \(\mathbf{C}\) , the image under the canonical injection \(\mathbf{Z} \to \mathbf{C}\) of a non-trivial ultrafilter on \(\mathbf{Z}\) . Show that there exist non-maximal prime ideals \(\mathfrak{p}\) of \(\mathbf{A}\) such that \(\mathfrak{F}_{\mathfrak{p}} = \mathfrak{U}\) . (For all \(x \in \mathfrak{a}(\mathfrak{U})\) and all \(n \in \mathbf{Z}\) , let \(\omega_x(n)\) be the order of \(x\) at the point \(n\) ( \(\omega_x(n) = 0\) if \(x(n) \neq 0\) ); consider the set \(\mathfrak{p}\) of \(x \in \mathfrak{a}(\mathfrak{U})\) such that \(\lim_{\mathfrak{u}} \omega_x = +\infty\) .)
\item
  With the notation of (d), let m be the maximal ideal \(\mathfrak{a}(\mathfrak{U})\) . Show that the domain \(A_{m}\) is not completely integrally closed (consider the meromorphic function \(1/(\sin\pi z)\) ).
\end{enumerate}

¶ 13. Let A be a local integral domain and K its field of fractions. Consider the two following properties:

\begin{enumerate}
\def\labelenumi{(\roman{enumi})}
\item
  A is Henselian (Chapter III, § 4, Exercise 3).
\item
  For every finite extension \(\mathbf{L}\) of \(\mathbf{K}\) , the integral closure \(\mathbf{A}'\) of \(\mathbf{A}\) in \(\mathbf{L}\) is a local ring.
\end{enumerate}

Show that (i) implies (ii) and that the converse is true if A is integrally closed. (To see that (i) implies (ii), note that A' is the union of a right directed family of finite A-algebras and use the definition of a Henselian ring. To see that (ii) implies (i) if A is integrally closed, show that, if P ∈ A{[}X{]} is an irreducible monic polynomial and f: A → k is the canonical homomorphism of A onto its residue field k, \(\bar{f}(P)\) is irreducible in K{[}X{]} and consider the field L = K{[}X{]}/(f).)

If condition (ii) is fulfilled, the integral closure A' of A in any extension of K is a local ring. Consider the case of complete Hausdorff local rings; if A is a complete Hausdorff Noetherian local ring, every A-algebra contained in A' and which is a finitely generated A-module is a complete Hausdorff local ring.

\begin{enumerate}
\def\labelenumi{\arabic{enumi}.}
\setcounter{enumi}{13}
\tightlist
\item
  Let A be a completely integrally closed domain and K its field of fractions. Show that for every extension L of K the integral closure A' of A in
\end{enumerate}

L is a completely integrally closed domain. (Reduce it to the case where L is a finite quasi-Galois extension of K of degree m: if \(x \in L\) is such that there exists \(d \in A'\) such that \(dx^{n} \in A'\) for all \(n \geqslant 0\) , show first that it may be assumed that \(d \in A\) , then deduce that, if \(c_{i}\) ( \(1 \leqslant i \leqslant m\) ) are the coefficients of the minimal polynomial of x over K, then \(d^{m}c_{i}^{n} \in A\) for all \(n \geqslant 0\) .)

\begin{enumerate}
\def\labelenumi{\arabic{enumi}.}
\setcounter{enumi}{14}
\tightlist
\item
  Let K be a commutative field of characteristic 0 containing all the roots of unity and such that there exist Galois extensions of K whose Galois group is not solvable (cf.~for example Algebra, Chapter V, Appendix I, no. 1, Proposition 2). Let \(\Omega\) be an algebraic closure of K and let \(\Omega_0\) be the set of elements \(z \in \Omega\) such that the Galois extension of K generated by \(z\) in \(\Omega\) admits a solvable Galois group.
\end{enumerate}

\begin{enumerate}
\def\labelenumi{(\alph{enumi})}
\item
  Show that \(\Omega_0\) is a subgroup of \(\Omega\) distinct from \(\Omega\) (cf.~Algebra, Chapter V, § 10, no. 4, Theorem 1).
\item
  Let A be the subring of \(\Omega[X]\) consisting of the polynomials P such that \(\mathrm{P}(0) \in \Omega_{0}\) . Show that A is not integrally closed but that every element of its field of fractions \(\Omega(X)\) a power of which belongs to A is itself in A.
\end{enumerate}

\begin{enumerate}
\def\labelenumi{\arabic{enumi}.}
\setcounter{enumi}{15}
\item
  Let B be a ring, A a subring of B and f the ideal of A the annihilator of the A-module B/A. Let U be the complement in \(\mathbf{X} = \operatorname{Spec}(\mathbf{A})\) of the closed set \(\mathbf{V}(\mathfrak{f})\) and let \(u = {}^{a}\rho\) , where \(\rho: A \to B\) is the canonical injection. Show that the restriction of u to \({}^{-1}(U)\) is a homeomorphism of \({}^{-1}(U)\) onto U (for every element \(g \in f\) , consider the spectrum of the ring \(A_{g}\) identified with the open set \(X_{g} \subset U\) ).
\item
  Let K be a field, B the polynomial ring \(K[X_{n}]_{n\in N}\) and A the subring of B generated by the monomials \(X_{n}^{2}\) and \(X_{n}^{3}\) for all \(n\in N\) ; show that B is the integral closure of A and that the conductor of B in A is reduced to 0; but, if \(S=A-\{0\}\) , the conductor of \(S^{-1}B\) in \(S^{-1}A\) is not reduced to 0 and \(S^{-1}A\) is integrally closed.
\item
  Let A be an integral domain, K its field of fractions, M a finitely generated A-module, u an endomorphism of M and \(u \otimes 1\) the corresponding endomorphism of the finite dimensional vector K-space \(M_{(K)} = M \otimes_{A} K\) . Show that the coefficients of the characteristic polynomial of \(u \otimes 1\) are integral over A (use condition \((E_{\mathrm{III}})\) of Theorem 1 of no. 1).
\item
  \begin{enumerate}
  \def\labelenumii{(\alph{enumii})}
  \tightlist
  \item
    Let A be an integrally closed domain, K its field of fractions, L a separable algebraic extension of K, x an element of L integral over A, \(K'\) the field \(\mathrm{K}(x) \subset \mathrm{L}\) and F the minimal polynomial of x over K. If x is of degree n over K, show that the integral closure of A in \(K'\) is contained in the A-module generated by the elements \(1/\mathrm{F}'(x)\) , \(x/\mathrm{F}'(x)\) , \(\ldots\) , \(x^{n-1}/\mathrm{F}'(x)\) of \(K'\) . (If \(z \in K'\) is integral over A, write
  \end{enumerate}
\end{enumerate}

\[
z \mathrm{F} ^ {\prime} (x) = g (x), \quad \text { where } \quad g (\mathbf {X}) = b _ {0} + b _ {1} \mathbf {X} + \dots + b _ {n - 1} \mathbf {X} ^ {n - 1} \in \mathbf {K} [ \mathbf {X} ].
\]

Show with the aid of Lagrange's interpolation formula that

\[
g (\mathbf {X}) = \operatorname{Tr} _ {\mathbf {K} ^ {\prime} [ \mathbf {X} ] / \mathbf {K} [ \mathbf {X} ]} \left(z. \frac {\mathrm{F} (\mathbf {X})}{\mathbf {X} - x}\right) \cdot)
\]

\begin{enumerate}
\def\labelenumi{(\alph{enumi})}
\setcounter{enumi}{1}
\tightlist
\item
  Let A be an integrally closed domain, K its field of fractions and p the characteristic exponent of K. Suppose that the subring \(A^{1/p}\) of \(K^{1/p}\) consisting of the p-th roots of the elements of A is a finitely generated A-module. Show that for every finite extension E of K the integral closure of A in E is a finitely generated A-module (reduce to the case where E is a quasi-Galois extension).
\end{enumerate}

¶ * 20. Let \(\mathbf{K}_0\) be a perfect field of characteristic \(p > 0\) , \(\mathbf{K}\) the field of rational functions \(\mathrm{K}_0(\mathrm{X}_n)_{n \in \mathbb{N}}\) and \(\mathbf{B}\) the ring of formal power series \(\mathrm{K}[[\mathrm{Z}]]\) , where \(\mathbf{Z}\) is an indeterminate; if \(m\) is the maximal ideal of \(B\), \(B\) is Hausdorff and complete with the \(m\) -adic topology and is a discrete valuation ring. Let \(L = K((Z))\) be the field of fractions of \(B\) and \(E_0\) the subfield of \(L\) generated by \(L^p\) and the elements \(Z\) and \(X_n\) (for \(n \in N\) ).

\begin{enumerate}
\def\labelenumi{(\alph{enumi})}
\item
  Show that the element \(c = \sum_{n=0}^{\infty} \mathbf{X}_n \mathbf{Z}^n\) of \(L\) does not belong to \(E_0\) .
\item
  Let E be a maximal element of the set of subfields of L containing \(E_{0}\) and not containing c; we write \(A = B \cap E\) . Show that A is a discrete valuation ring (and therefore Noetherian) whose field of fractions is E; if \(m'\) is the maximal ideal of A, then \(m'^{k} = m^{k} \cap A\) and A is dense in B with the m-adic topology; \(L = K(c)\) is an extension of K of degree p and B is the integral closure of A in L but B is not a finitely generated A-module (use Chapter III, §3, Exercise 9 (b)). *
\end{enumerate}

\begin{enumerate}
\def\labelenumi{\arabic{enumi}.}
\setcounter{enumi}{20}
\tightlist
\item
  \begin{enumerate}
  \def\labelenumii{(\alph{enumii})}
  \tightlist
  \item
    Let \(K_{0}\) be a perfect field of characteristic p \textgreater{} 0, L the field of rational functions \(\mathrm{K}_{0}(\mathrm{X}_{n})_{n\in\mathbf{N}}\) , \(A'\) the ring of formal power series \(L[[Y,Z]]\) , where Y and Z are two indeterminates, and A the tensor product ring \(K[[Y,Z]]\otimes_{K}L\) , where we write \(K = L^{p} \subset L\) ; A is a Noetherian local ring, which is not complete and whose completion is identified with \(A'\) (Chapter III, §3, Exercise 17) and \(A'\) is integral over A. Let \(c_{n}\) denote the element
  \end{enumerate}
\end{enumerate}

\[
\mathrm{X} _ {n} \mathrm{Y} + \mathrm{X} _ {n + 1} \mathrm{YZ} + \mathrm{X} _ {n + 2} \mathrm{YZ} ^ {2} + \dots
\]

of A'. Let B be the subring of A' generated by A and the elements \(c_{n}\) ( \(n \geqslant 0\) ). If C = A{[}c₀{]}, C is Noetherian and the ring B (which is not integrally closed) is contained in the integral closure of C and contains C; but show that B is not Noetherian (observe that \(c_{n}\) does not belong to the ideal of B generated by the \(c_{i}\) of index i \textless{} n).

\begin{enumerate}
\def\labelenumi{(\alph{enumi})}
\setcounter{enumi}{1}
\tightlist
\item
  Suppose \(p = 2\) ; \(K_0\) , \(K\) and \(L\) being used with the same meaning as in (a), consider this time the ring of formal power series \(A' = L[[Y, Z, T]]\) in three indeterminates, the subring
\end{enumerate}

\[
\mathrm{A} = \mathrm{K} [ [ \mathrm{Y}, \mathrm{Z}, \mathrm{T} ] ] \otimes_ {\mathrm{K}} \mathrm{L},
\]

and write

\[
b = \mathrm{Y} \sum_ {n \geqslant 0} \mathrm{X} _ {2 n} \mathrm{T} ^ {n} + \mathrm{Z} \sum_ {n \geqslant 0} \mathrm{X} _ {2 n + 1} \mathrm{T} ^ {n}.
\]

Show that the ring A{[}b{]} is Noetherian but that its integral closure B is not a Noetherian ring. (As \(b^{2} \in A\) , every element of the field of fractions of A{[}b{]} is of the form \(P + Qb\) , where P and Q are linear combinations of a finite number of \(X_{k}\) with coefficients in the field of fractions of K{[}{[}Y, Z, T{]}{]}; for such an element to be integral over A{[}b{]}, show that it is necessary and sufficient that it belong to A'. Show that B is generated by the elements

\[
b _ {i} = \mathrm{Y} \sum_ {n \geqslant 0} \mathrm{X} _ {2 (i + n)} \mathrm{T} ^ {n} + \mathrm{Z} \sum_ {n \geqslant 0} \mathrm{X} _ {2 (i + n) + 1} \mathrm{T} ^ {n} \quad \text { for } \quad i \geqslant 0
\]

and conclude the argument as in (a).)

\begin{enumerate}
\def\labelenumi{\arabic{enumi}.}
\setcounter{enumi}{21}
\item
  Let A be an integrally closed domain and K its field of fractions; show that for every extension L of K there exists an integrally closed subdomain B of L whose field of fractions is L and such that \(B \cap K = A\) (consider L as an algebraic extension of a pure transcendental extension \(E = K(X_t)_{t \in I}\) of K and thus reduce it to the case where L is an algebraic extension of K, using Exercise 11 (a)).
\item
  \begin{enumerate}
  \def\labelenumii{(\alph{enumii})}
  \tightlist
  \item
    Let K be a commutative field, n an integer prime to the characteristic of K and P a polynomial in K{[}X{]} with only simple roots (in an algebraic closure of K). If we write \(f(\mathbf{X}, \mathbf{Y}) = \mathbf{Y}^{n} - \mathbf{P}(\mathbf{X})\) , show that, if K contains the n-th roots of unity, the domain \(\mathbf{A} = \mathbf{K}[\mathbf{X}, \mathbf{Y}] / (f)\) is integrally closed and has in its field of fractions the separable extension L of K(X) generated by a root y of f (considered as a polynomial in K(X){[}Y{]}). (Show that, if an element z of L is integral over K{[}X{]}, it belongs to A; for this, write z in the form
  \end{enumerate}
\end{enumerate}

\[
z = \sum_ {i = 0} ^ {n - 1} a _ {i} (\mathbf {X}) y ^ {i},
\]

where \(a_{i}(\mathbf{X}) \in \mathbf{K}(\mathbf{X})\) . Show that the elements \(a_{i}(\mathbf{X})y^{i} (0 \leqslant i \leqslant n - 1)\) are integral over \(\mathbf{K}[\mathbf{X}]\) and deduce that the \((a_{i}(\mathbf{X}))^{n}(\mathbf{P}(\mathbf{X}))^{i}\) are elements of \(\mathbf{K}[\mathbf{X}]\) ; conclude that so are the \(a_{i}(\mathbf{X})\) .

\begin{enumerate}
\def\labelenumi{(\alph{enumi})}
\setcounter{enumi}{1}
\tightlist
\item
  Suppose that K is an imperfect field of characteristic \(p \neq 2\) ; let \(a \in K\) be an element not belonging to \(K^{p}\) , let E be the field \(\mathrm{K}(a^{1/p})\) and let us write \(\mathrm{P}(\mathrm{X}) = \mathrm{X}^{p} - a, f(\mathrm{X}, \mathrm{Y}) = \mathrm{Y}^{2} - \mathrm{P}(\mathrm{X})\) . If A is the integrally closed domain \(\mathrm{K}[\mathrm{X}, \mathrm{Y}]/(f)\) , show that \(B = E \otimes_{K} A\) is an integral domain but not integrally closed (consider the element \(y/(\mathrm{X} - a^{1/p})\) of the field of fractions of B).
\end{enumerate}

\begin{enumerate}
\def\labelenumi{\arabic{enumi}.}
\setcounter{enumi}{23}
\tightlist
\item
  Let \(\Delta\) be a torsion-free commutative group written additively. Let A be a commutative ring, B a commutative A-algebra and \(h\colon\mathbf{A}\to\mathbf{B}\) the ring homomorphism defining the algebra structure on B. Then \(h^{(\Delta)}\colon\mathbf{A}^{(\Delta)}\to\mathbf{B}^{(\Delta)}\) (Algebra,
\end{enumerate}

Chapter II, § 11, Exercise 1) is a (ring) homomorphism of the algebra \(\mathbf{A}^{(\Delta)}\) of the group \(\Delta\) with respect to \(\mathbf{A}\) to the algebra \(\mathbf{B}^{(\Delta)}\) of the group \(\Delta\) with respect to \(\mathbf{B}\) . Let \(\mathbf{P} = \sum_{\alpha \in \Delta} b_{\alpha} \otimes \mathbf{X}^{\alpha}\) be an element of \(\mathbf{B}^{(\Delta)}\) . Show that for \(\mathbf{P}\) to be integral over \(\mathbf{A}^{(\Delta)}\) , it is necessary and sufficient that each of the elements \(b_{\alpha} \in \mathbf{B}\) is integral over \(\mathbf{A}\) . (Reduce it to the case where \(\Delta\) is finitely generated and then to the case \(\Delta = \mathbf{Z}\) .) Show that if \(\mathbf{A}\) is an integral domain and integrally closed so is \(\mathbf{A}^{(\Delta)}\) .

\begin{enumerate}
\def\labelenumi{\arabic{enumi}.}
\setcounter{enumi}{24}
\tightlist
\item
  Let \(\Delta\) be a totally ordered commutative group. Extend Propositions 20 and 21 of no. 8 to graded rings of type \(\Delta\) (for Proposition 20 use Exercise 24; for Proposition 21 generalize Lemma 4 of no. 8 to the case where \(\Delta\) is finitely generated and then prove Proposition 21 by taking the direct limit).
\end{enumerate}

¶ 26. (a) Let A be an integral domain such that for all \(a \neq 0\) in A, the ideal \(Aa\) is the intersection of a family \((q_t)\) of saturations of \(Aa\) with respect to prime ideals \(p_t\) which are minimal among those which contain \(a\) , such that the domains \(A_{p_t}\) are integrally closed (resp. completely integrally closed). Then show that A is integrally closed (resp. completely integrally closed).

\begin{enumerate}
\def\labelenumi{(\alph{enumi})}
\setcounter{enumi}{1}
\item
  Let A be a strongly Laskerian local ring (Chapter IV, § 2, Exercise 28) and m its maximal ideal. Suppose that, for some \(a \in \mathbf{A}\) , m is an immersed prime ideal weakly associated with A/Aa (Chapter IV, § 1, Exercise 17). Show that Aa: m is contained in every ideal which is primary for every prime ideal p ≠ m weakly associated with A/Aa (consider the transporter q: m for such a primary ideal q, cf.~Chapter IV, § 2, Exercise 30). If further A is assumed to be an integral domain, show that the relation p.~(Aa: m) = Aa, for a prime ideal p minimal among those which contain Aa, is impossible (if q is the saturation of A with respect to p, note that the above relation would imply that qA \(_{p}\) = qpA \(_{p}\) in A \(_{p}\) and conclude with the aid of Exercise 29 (d) of Chapter IV, § 2).
\item
  Show that, if A is a strongly Laskerian integral domain and there exists an element \(a \neq 0\) of A such that there are prime ideals which are weakly associated with A/Aa and immersed, A is not completely integrally closed. (Reduce it to the case considered in (b); in the notation of (b), there then exists \(b \in \mathbf{A}a\) : m such that \(b \notin \mathbf{A}a\) and \(bp \subset am\) ; deduce by induction on \(n\) that \(b^n p \subset a^n m\) for all \(n\) .)
\item
  Deduce in particular from (a) and (c) that, for a Noetherian integral domain to be integrally closed, it is necessary and sufficient that for all \(a \neq 0\) in A the prime ideals \(\mathfrak{p}\) associated with A/Aa be not immersed and be such that \(\mathrm{A}_{\mathfrak{p}}\) is integrally closed (cf.~Chapter VII, § 1, no. 4, Proposition 8).
\end{enumerate}

\begin{enumerate}
\def\labelenumi{\arabic{enumi}.}
\setcounter{enumi}{26}
\tightlist
\item
  Let A be an integrally closed but not completely integrally closed domain and K its field of fractions; let \(a \neq 0\) be a non-invertible element of A for which there exists \(d \neq 0\) in A such that \(da^{-n} \in A\) for every integer \(n \geqslant 0\) .
\end{enumerate}

Show that in the field of formal power series \(\mathbf{K}((\mathbf{X}))\) there exists a formal power series \(\sum_{k=1}^{\infty} u_k \mathbf{X}^k = f(\mathbf{X})\) satisfying the equation

\[
(f (\mathbf {X})) ^ {2} - a f (\mathbf {X}) + \mathbf {X} = 0,
\]

hence integral over A{[}{[}X{]}{]}, belonging to the field of fractions of A{[}{[}X{]}{]} but not belonging to A{[}{[}X{]}{]}, so that A{[}{[}X{]}{]} is not integrally closed.

\begin{enumerate}
\def\labelenumi{\arabic{enumi}.}
\setcounter{enumi}{27}
\item
  Let \(k\) be a field and \(\mathbf{A}_0\) the polynomial ring \(k[\mathrm{X}_n, \mathrm{Y}_n]_{n \in \mathbf{Z}}\) in two infinite families of indeterminates. Let \(G\) be an infinite monogenous group (therefore isomorphic to \(\mathbf{Z}\) ) and let \(\sigma\) be a generator of \(\mathcal{G}\) . \(G\) is defined as operating on \(\mathbf{A}_0\) by the conditions \(\sigma(\mathrm{Y}_n) = \mathrm{Y}_n\) and \(\sigma(\mathrm{X}_n) = \mathrm{X}_{n+1}\) for all \(n \in \mathbf{Z}\) . Let \(\Im\) be the ideal of \(\mathbf{A}_0\) generated by the elements \(\mathrm{Y}_n(\mathrm{X}_n - \mathrm{X}_{n+1})\) for all \(n \in \mathbf{Z}\) and let \(A\) be the quotient ring \(A_0 / \Im\) ; as \(\sigma(\Im) \subset \Im\) , the group \(\mathcal{G}\) also operates on \(A\) by passing to the quotient. Let \(S\) be the multiplicative subset of \(A\) generated by the canonical images of the \(Y_n\) in \(A\) , so that \(S\) consists of the elements invariant under \(G\) . Show that \(S^{-1}(A^{\mathcal{G}}) \neq (S^{-1}A)^{\mathcal{G}}\) .
\item
  Let k be a field of characteristic \(\neq2\) , A the polynomial ring k{[}T{]} in one indeterminate and a, b two elements of k such that \(a\neq0\) , \(b\neq0\) and \(a\neq b\) . Let \(\mathbf{K}=k(\mathbf{T})\) be the field of fractions of A and L, M the extensions of K adjoining to K respectively a root of \(\mathbf{X}^{2}-\mathbf{T}(\mathbf{T}-a)\) and a root of \(\mathbf{X}^{2}-\mathbf{T}(\mathbf{T}-b)\) . Let B, C be the integral closures of A in L and M respectively, which are finitely generated free A-modules. Show that \(L\otimes_{K}M\) is a field, a separable finite extension of K, but that \(B\otimes_{A}C\) , which is canonically identified with a subring of \(L\otimes_{K}M\) , is not the integral closure of A in \(L\otimes_{K}M\) .
\end{enumerate}
